\documentclass[12pt, a4paper]{article}

% ==========================================
% 1. COMPILER, LANGUAGE & FONT SETUP
% ==========================================
\usepackage{fontspec}
\usepackage{polyglossia}
\usepackage[protrusion=false]{microtype}
\setmainlanguage{bengali}
\setotherlanguage{english}
\newfontfamily\bengalifont[Script=Bengali, AutoFakeBold=2.5, AutoFakeSlant=0.2]{Kalpurush}

% ==========================================
% 2. MATHEMATICS, UNITS & FORMATTING
% ==========================================
\usepackage{amsmath, amssymb, siunitx}
\usepackage[margin=1in]{geometry}
\usepackage{setspace}
\usepackage{enumitem}
\usepackage{microtype} % For better typography
\onehalfspacing % Better line spacing for readability

% ==========================================
% 3. COLORS & BEAUTIFICATION PACKAGES
% ==========================================
\usepackage[dvipsnames, table]{xcolor}
\usepackage[skins, breakable, theorems]{tcolorbox}
\usepackage{tikz}
\usetikzlibrary{shadows, arrows.meta, positioning, decorations.pathmorphing, calc, intersections, angles, quotes}
\usepackage{enumitem}
\setlist[itemize,1]{label=\textendash}

% Define custom beautiful colors
\definecolor{primary}{HTML}{0056b3}
\definecolor{secondary}{HTML}{e7f1ff}
\definecolor{success}{HTML}{198754}
\definecolor{successbg}{HTML}{d1e7dd}
\definecolor{accent}{HTML}{fd7e14}
\definecolor{darkgray}{HTML}{343a40}

% Custom Tcolorbox Styles
\tcbset{
    questionbox/.style={
        enhanced, colback=secondary, colframe=primary, arc=2mm, boxrule=1pt,
        fonttitle=\bfseries\Large, title=#1,
        drop fuzzy shadow=black!10,
        attach boxed title to top left={xshift=5mm, yshift=-3mm},
        boxed title style={colback=primary, arc=1mm}
    },
    answerbox/.style={
        enhanced, colback=successbg, colframe=success, arc=1.5mm, boxrule=1pt,
        left=2mm, right=2mm, top=2mm, bottom=2mm,
        drop fuzzy shadow=black!10
    },
    solutionbox/.style={
        enhanced, colback=white, colframe=darkgray, arc=2mm, boxrule=1pt,
        fonttitle=\bfseries\large, title=#1, breakable,
        coltitle=white, colbacktitle=darkgray,
        drop fuzzy shadow=black!10,
        attach boxed title to top center={yshift=-3mm},
        boxed title style={arc=1mm}
    }
}

\begin{document}

% ------------------------------------------
% QUESTION SECTION
% ------------------------------------------
\begin{tcolorbox}[questionbox={প্রশ্ন (Question) 1}]
    \noindent
    আর্গন্ড সমতলে একটি চলমান জটিল সংখ্যা $z$ এমনভাবে অবস্থান করে যেন $\text{Re}\left(\frac{z-1}{z+1}\right) = 2$ হয়, তবে $z$ এর সঞ্চারপথটি জ্যামিতিকভাবে কী নির্দেশ করে?
    
    \vspace{0.8em}
    \begin{english}
    \noindent\textit{\color{darkgray}
    A point $z$ moves in the Argand plane such that $\text{Re}\left(\frac{z-1}{z+1}\right) = 2$. What does the locus of $z$ geometrically represent?
    }
    \end{english}
\end{tcolorbox}

% ------------------------------------------
% OPTIONS SECTION
% ------------------------------------------
\begin{itemize}[label={}, leftmargin=1cm, itemsep=0.5em]
    \item[\textbf{A)}] $y$-অক্ষের সমান্তরাল একটি সরলরেখা (A straight line parallel to the $y$-axis)
    \item[\textbf{B)}] $(-2, 0)$ কেন্দ্রে অবস্থিত $1$ একক ব্যাসার্ধের একটি বৃত্ত (A circle centered at $(-2, 0)$ with a radius of $1$ unit)
    \item[\textbf{C)}] $(-2, 0)$ কেন্দ্রে অবস্থিত $2$ একক ব্যাসার্ধের একটি বৃত্ত (A circle centered at $(-2, 0)$ with a radius of $2$ units)
    \item[\textbf{D)}] একটি উপবৃত্ত যার বৃহৎ অক্ষ $x$-অক্ষ বরাবর (An ellipse with its major axis along the $x$-axis)
\end{itemize}

% ------------------------------------------
% ANSWER HIGHLIGHT
% ------------------------------------------
\vspace{0.5em}
\begin{tcolorbox}[answerbox]
    \textbf{\color{success} উত্তর:} B) $(-2, 0)$ কেন্দ্রে অবস্থিত $1$ একক ব্যাসার্ধের একটি বৃত্ত
\end{tcolorbox}
\vspace{1em}

% ------------------------------------------
% SOLUTION SECTION
% ------------------------------------------
\begin{tcolorbox}[solutionbox={সমাধান (Solution)}]
    \noindent
    ধরি, $z = x + iy$। প্রদত্ত রাশিটি হলো:
    \[ \frac{z-1}{z+1} = \frac{(x-1) + iy}{(x+1) + iy} \]
    লব ও হরকে হর-এর অনুবন্ধী দ্বারা গুণ করে পাই:
    \begin{align*}
        \frac{z-1}{z+1} &= \frac{((x-1) + iy)((x+1) - iy)}{(x+1)^2 + y^2} \\
        &= \frac{(x^2 - 1 + y^2) + i(y(x+1) - y(x-1))}{(x+1)^2 + y^2} \\
        &= \frac{(x^2 + y^2 - 1) + 2iy}{(x+1)^2 + y^2}
    \end{align*}
    এখন, বাস্তব অংশ সমীকৃত করে পাই:
    \begin{align*}
        \text{Re}\left(\frac{z-1}{z+1}\right) &= \frac{x^2 + y^2 - 1}{(x+1)^2 + y^2} = 2 \\
        \Rightarrow x^2 + y^2 - 1 &= 2(x^2 + 2x + 1 + y^2) \\
        \Rightarrow x^2 + y^2 - 1 &= 2x^2 + 4x + 2 + 2y^2 \\
        \Rightarrow x^2 + y^2 + 4x + 3 &= 0 \\
        \Rightarrow (x+2)^2 + y^2 &= 1
    \end{align*}
    এটি একটি বৃত্তের সমীকরণ নির্দেশ করে যার কেন্দ্র $(-2, 0)$ এবং ব্যাসার্ধ $r = \sqrt{1} = 1$ একক।
\end{tcolorbox}

\vspace{1em}

% ------------------------------------------
% QUESTION SECTION
% ------------------------------------------
\begin{tcolorbox}[questionbox={প্রশ্ন (Question) 2}]
    \noindent
    যদি $a \neq 1$ এমন একটি জটিল সংখ্যা হয় যার জন্য $\frac{a-1}{a+1}$ একটি সম্পূর্ণ কাল্পনিক সংখ্যা হয়, তবে $a + \frac{1}{a}$ এর মান কেমন হবে?
    
    \vspace{0.8em}
    \begin{english}
    \noindent\textit{\color{darkgray}
    If $a \neq 1$ is a complex number such that $\frac{a-1}{a+1}$ is a purely imaginary number, then the value of $a + \frac{1}{a}$ is:
    }
    \end{english}
\end{tcolorbox}

% ------------------------------------------
% OPTIONS SECTION
% ------------------------------------------
\begin{itemize}[label={}, leftmargin=1cm, itemsep=0.5em]
    \item[\textbf{A)}] যেকোনো বাস্তব সংখ্যা (A real number)
    \item[\textbf{B)}] একটি সম্পূর্ণ কাল্পনিক সংখ্যা (A purely imaginary number)
    \item[\textbf{C)}] সর্বদা $0$ এর সমান (Always equal to $0$)
    \item[\textbf{D)}] $1$ এর চেয়ে বড় যেকোনো ধনাত্মক বাস্তব সংখ্যা (Any positive real number greater than $1$)
\end{itemize}

% ------------------------------------------
% ANSWER HIGHLIGHT
% ------------------------------------------
\vspace{0.5em}
\begin{tcolorbox}[answerbox]
    \textbf{\color{success} উত্তর:} A) যেকোনো বাস্তব সংখ্যা
\end{tcolorbox}
\vspace{1em}

% ------------------------------------------
% SOLUTION SECTION
% ------------------------------------------
\begin{tcolorbox}[solutionbox={সমাধান (Solution)}]
    \noindent
    ধরি, $a = x + iy$। আমরা পাই:
    \[ \frac{a-1}{a+1} = \frac{(x-1) + iy}{(x+1) + iy} = \frac{(x^2 + y^2 - 1) + 2iy}{(x+1)^2 + y^2} \]
    যেহেতু এটি সম্পূর্ণ কাল্পনিক সংখ্যা, তাই এর বাস্তব অংশ অবশ্যই $0$ হবে:
    \[ \text{Re}\left(\frac{a-1}{a+1}\right) = 0 \Rightarrow x^2 + y^2 - 1 = 0 \Rightarrow |a|^2 = 1 \Rightarrow |a| = 1 \]
    এখন,
    \[ a + \frac{1}{a} = a + \frac{\bar{a}}{|a|^2} \]
    যেহেতু $|a| = 1$, সেহেতু:
    \[ a + \frac{1}{a} = a + \bar{a} \]
    আমরা জানি $a + \bar{a} = 2x$ (যা একটি বাস্তব সংখ্যা)। 
    সুতরাং, $a + \frac{1}{a}$ সর্বদা একটি বাস্তব সংখ্যা হবে।
\end{tcolorbox}


% ------------------------------------------
% QUESTION SECTION
% ------------------------------------------
\begin{tcolorbox}[questionbox={প্রশ্ন (Question) 3}]
    \noindent
    আর্গন্ড সমতলে জটিল সংখ্যা $z$, $iz$ এবং $z + iz$ শীর্ষবিন্দু বিশিষ্ট ত্রিভুজের ক্ষেত্রফল $50$ বর্গ একক হলে, জটিল সংখ্যা $z$ এর পরমমান $|z|$ এর মান কত?
    
    \vspace{0.8em}
    \begin{english}
    \noindent\textit{\color{darkgray}
    If the area of the triangle formed by the points $z$, $iz$, and $z + iz$ in the Argand plane is $50$ square units, then what is the value of the modulus $|z|$ of the complex number $z$?
    }
    \end{english}
\end{tcolorbox}

% ------------------------------------------
% OPTIONS SECTION
% ------------------------------------------
\begin{itemize}[label={}, leftmargin=1cm, itemsep=0.5em]
    \item[\textbf{A)}] $5$
    \item[\textbf{B)}] $10$
    \item[\textbf{C)}] $15$
    \item[\textbf{D)}] $20$
\end{itemize}

% ------------------------------------------
% ANSWER HIGHLIGHT
% ------------------------------------------
\vspace{0.5em}
\begin{tcolorbox}[answerbox]
    \textbf{\color{success} উত্তর:} B) $10$
\end{tcolorbox}
\vspace{1em}

% ------------------------------------------
% SOLUTION SECTION
% ------------------------------------------
\begin{tcolorbox}[solutionbox={সমাধান (Solution)}]
    \noindent
    ত্রিভুজের শীর্ষবিন্দু তিনটি হলো $A(z)$, $B(iz)$ এবং বা $C(z+iz)$। বাহুগুলোর দৈর্ঘ্য পরিমাপ করে পাই:
    \begin{align*}
        AB &= |iz - z| = |z(i - 1)| = |z||i-1| = |z|\sqrt{1^2+(-1)^2} = |z|\sqrt{2} \\
        BC &= |(z+iz) - iz| = |z| \\
        CA &= |(z+iz) - z| = |iz| = |z|
    \end{align*}
    দেখা যাচ্ছে যে, বাহুর দৈর্ঘ্য $BC = CA = |z|$ এবং $AB = |z|\sqrt{2}$। যেহেতু $BC^2 + CA^2 = |z|^2 + |z|^2 = 2|z|^2 = AB^2$, তাই ত্রিভুজটি একটি সমকোণী সমদ্বিবাহু ত্রিভুজ, যার অতিভুজ $AB$ এবং অপর দুটি লম্ব বাহুর দৈর্ঘ্য $|z|$।
    
    ত্রিভুজটির ক্ষেত্রফল:
    \[ \text{Area} = \frac{1}{2} \times \text{ভূমি} \times \text{উচ্চতা} = \frac{1}{2} |z|^2 \]
    শর্তানুসারে:
    \[ \frac{1}{2} |z|^2 = 50 \Rightarrow |z|^2 = 100 \Rightarrow |z| = 10 \]
\end{tcolorbox}

\vspace{1em}

% ------------------------------------------
% QUESTION SECTION
% ------------------------------------------
\begin{tcolorbox}[questionbox={প্রশ্ন (Question) 4}]
    \noindent
    জটিল সমীকরণের বৈশিষ্ট্য বিশ্লেষণ করে $iz^3 + z^2 - z + i = 0$ সমীকরণটির জন্য $|z|$ এর সঠিক মান নির্ণয় করো।
    
    \vspace{0.8em}
    \begin{english}
    \noindent\textit{\color{darkgray}
    Analyzing the properties of complex equations, determine the correct value of $|z|$ for the equation $iz^3 + z^2 - z + i = 0$.
    }
    \end{english}
\end{tcolorbox}

% ------------------------------------------
% OPTIONS SECTION
% ------------------------------------------
\begin{itemize}[label={}, leftmargin=1cm, itemsep=0.5em]
    \item[\textbf{A)}] $4$
    \item[\textbf{B)}] $2$
    \item[\textbf{C)}] $3$
    \item[\textbf{D)}] $1$
\end{itemize}

% ------------------------------------------
% ANSWER HIGHLIGHT
% ------------------------------------------
\vspace{0.5em}
\begin{tcolorbox}[answerbox]
    \textbf{\color{success} উত্তর:} D) $1$
\end{tcolorbox}
\vspace{1em}

% ------------------------------------------
% SOLUTION SECTION
% ------------------------------------------
\begin{tcolorbox}[solutionbox={সমাধান (Solution)}]
    \noindent
    প্রদত্ত সমীকরণটি হলো: $iz^3 + z^2 - z + i = 0$
    
    সমীকরণটিকে উৎপাদকে বিশ্লেষণ করার চেষ্টা করি:
    \begin{align*}
        iz^2(z - i) - 1(z - i) &= 0 \\
        \Rightarrow (iz^2 - 1)(z - i) &= 0
    \end{align*}
    অতএব, দুটি সম্ভাব্য ক্ষেত্র তৈরি হয়:
    \begin{itemize}
        \item \textbf{ক্ষেত্র-১:} $z - i = 0 \Rightarrow z = i \Rightarrow |z| = |i| = 1$
        \item \textbf{ক্ষেত্র-২:} $iz^2 - 1 = 0 \Rightarrow iz^2 = 1 \Rightarrow z^2 = \frac{1}{i} = -i$
    \end{itemize}
    উভয়পক্ষে মডুলাস নিয়ে পাই:
    \[ |z^2| = |-i| \Rightarrow |z|^2 = 1 \Rightarrow |z| = 1 \]
    উভয় ক্ষেত্রেই জটিল সংখ্যা $z$ এর পরমমান $|z| = 1$ একক।
\end{tcolorbox}
% ------------------------------------------
% QUESTION SECTION
% ------------------------------------------
\begin{tcolorbox}[questionbox={প্রশ্ন (Question) 5}]
    \noindent
    আর্গন্ড সমতলে একটি জটিল সংখ্যা $z = 3+4i$ কে অন্য একটি স্থির বিন্দু $z_0 = 1+i$ এর সাপেক্ষে ঘড়ির কাঁটার বিপরীত দিকে $90^\circ$ কোণে আবর্তন করানো হলে উৎপন্ন নতুন জটিল সংখ্যাটি নিচের কোনটি হবে?
    
    \vspace{0.8em}
    \begin{english}
    \noindent\textit{\color{darkgray}
    In the Argand plane, if a complex number $z = 3+4i$ is rotated about a fixed point $z_0 = 1+i$ through an angle of $90^\circ$ in the counter-clockwise direction, then the new complex number is:
    }
    \end{english}
\end{tcolorbox}

% ------------------------------------------
% OPTIONS SECTION
% ------------------------------------------
\begin{itemize}[label={}, leftmargin=1cm, itemsep=0.5em]
    \item[\textbf{A)}] $-2 + 3i$
    \item[\textbf{B)}] $-2 + 4i$
    \item[\textbf{C)}] $2 - 3i$
    \item[\textbf{D)}] $-1 + 3i$
\end{itemize}

% ------------------------------------------
% ANSWER HIGHLIGHT
% ------------------------------------------
\vspace{0.5em}
\begin{tcolorbox}[answerbox]
    \textbf{\color{success} উত্তর:} A) $-2 + 3i$
\end{tcolorbox}
\vspace{1em}

% ------------------------------------------
% SOLUTION SECTION
% ------------------------------------------
\begin{tcolorbox}[solutionbox={সমাধান (Solution)}]
    \noindent
    আমরা জানি, জটিল সমতলে যেকোনো জটিল সংখ্যা $z$ কে অন্য একটি স্থির বিন্দু $z_0$ এর সাপেক্ষে ঘড়ির কাঁটার বিপরীত দিকে $\theta$ কোণে আবর্তন করালে নতুন জটিল সংখ্যা $z'$ পাওয়া যায়, যা নিম্নোক্ত সমীকরণটি সিদ্ধ করে:
    \[ z' - z_0 = (z - z_0) e^{i\theta} \]
    এখানে প্রদত্ত মানগুলো হলো:
    \begin{itemize}
        \item $z = 3+4i$
        \item $z_0 = 1+i$
        \item $\theta = 90^\circ \Rightarrow e^{i\theta} = e^{i\pi/2} = i$
    \end{itemize}
    প্রথমে $z - z_0$ নির্ণয় করি:
    \[ z - z_0 = (3+4i) - (1+i) = 2 + 3i \]
    এখন একে $i$ দ্বারা গুণ করে পাই:
    \[ i(z - z_0) = i(2 + 3i) = 2i + 3i^2 = -3 + 2i \quad (\text{যেহেতু } i^2 = -1) \]
    তাহলে নতুন জটিল সংখ্যাটি হবে:
    \begin{align*}
        z' &= z_0 + i(z - z_0) = (1+i) + (-3+2i) \\
        \Rightarrow z' &= (1 - 3) + i(1 + 2) = -2 + 3i
    \end{align*}
    সুতরাং, নতুন জটিল সংখ্যাটি হবে $-2 + 3i$।
\end{tcolorbox}

\vspace{1em}

% ------------------------------------------
% QUESTION SECTION
% ------------------------------------------
\begin{tcolorbox}[questionbox={প্রশ্ন (Question) 6}]
    \noindent
    যদি $z$ একটি জটিল সংখ্যা হয় যা $|z+1| = z + 2(1+i)$ সমীকরণটি সিদ্ধ করে, তবে $z$ এর মান নিচের কোনটি?
    
    \vspace{0.8em}
    \begin{english}
    \noindent\textit{\color{darkgray}
    If $z$ is a complex number satisfying the equation $|z+1| = z + 2(1+i)$, then $z$ is:
    }
    \end{english}
\end{tcolorbox}

% ------------------------------------------
% OPTIONS SECTION
% ------------------------------------------
\begin{itemize}[label={}, leftmargin=1cm, itemsep=0.5em]
    \item[\textbf{A)}] $\frac{1}{2}(1 + 4i)$
    \item[\textbf{B)}] $\frac{1}{2}(3 + 4i)$
    \item[\textbf{C)}] $\frac{1}{2}(1 - 4i)$
    \item[\textbf{D)}] $\frac{1}{2}(3 - 4i)$
\end{itemize}

% ------------------------------------------
% ANSWER HIGHLIGHT
% ------------------------------------------
\vspace{0.5em}
\begin{tcolorbox}[answerbox]
    \textbf{\color{success} উত্তর:} C) $\frac{1}{2}(1 - 4i)$
\end{tcolorbox}
\vspace{1em}

% ------------------------------------------
% SOLUTION SECTION
% ------------------------------------------
\begin{tcolorbox}[solutionbox={সমাধান (Solution)}]
    \noindent
    ধরি, $z = x + iy$, যেখানে $x, y \in \mathbb{R}$। প্রদত্ত সমীকরণটি হলো:
    \begin{align*}
        |z+1| &= z + 2(1+i) \\
        \Rightarrow |x + iy + 1| &= (x + iy) + 2 + 2i \\
        \Rightarrow |(x+1) + iy| &= (x + 2) + i(y + 2) \quad \text{--- (1)}
    \end{align*}
    যেহেতু বামপক্ষের রাশিটি একটি মডুলাস বা পরমমান, তাই এটি অবশ্যই একটি বাস্তব সংখ্যা (যার কাল্পনিক অংশ শূন্য)। অতএব, ডানপক্ষের সংখ্যাটির কাল্পনিক অংশও অবশ্যই শূন্য হতে হবে। 
    
    কাল্পনিক অংশ সমীকৃত করে পাই:
    \[ y + 2 = 0 \Rightarrow y = -2 \]
    বাস্তব অংশ সমীকৃত করে পাই:
    \[ |(x+1) + iy| = x + 2 \]
    এখানে $y = -2$ বসিয়ে পাই:
    \begin{align*}
        |(x+1) - 2i| &= x + 2 \\
        \Rightarrow \sqrt{(x+1)^2 + (-2)^2} &= x + 2 \\
        \Rightarrow \sqrt{x^2 + 2x + 1 + 4} &= x + 2 \\
        \Rightarrow \sqrt{x^2 + 2x + 5} &= x + 2
    \end{align*}
    উভয় পক্ষকে বর্গ করে পাই (যেখানে $x + 2 \geq 0 \Rightarrow x \geq -2$):
    \begin{align*}
        x^2 + 2x + 5 &= (x + 2)^2 \\
        \Rightarrow x^2 + 2x + 5 &= x^2 + 4x + 4 \\
        \Rightarrow 2x + 5 &= 4x + 4 \\
        \Rightarrow 2x &= 1 \Rightarrow x = \frac{1}{2} \quad (\text{যা } x \geq -2 \text{ শর্তটি সিদ্ধ করে})
    \end{align*}
    অতএব, নির্ণেয় জটিল সংখ্যাটি হলো:
    \[ z = x + iy = \frac{1}{2} - 2i = \frac{1}{2}(1 - 4i) \]
\end{tcolorbox}
% ------------------------------------------
% QUESTION SECTION
% ------------------------------------------
\begin{tcolorbox}[questionbox={প্রশ্ন (Question) 7}]
    \noindent
    যদি জটিল সমীকরণ $z^2 + pz + q = 0$ এর মূলদ্বয় $z_1, z_2$ এমন হয় যে $|z_1| = |z_2| = 1$ (যেখানে $p, q$ বাস্তব সংখ্যা), তবে নিচের কোনটি সঠিক?
    
    \vspace{0.8em}
    \begin{english}
    \noindent\textit{\color{darkgray}
    If the roots $z_1, z_2$ of the complex quadratic equation $z^2 + pz + q = 0$ satisfy $|z_1| = |z_2| = 1$ (where $p, q$ are real numbers), then which of the following is correct?
    }
    \end{english}
\end{tcolorbox}

% ------------------------------------------
% OPTIONS SECTION
% ------------------------------------------
\begin{itemize}[label={}, leftmargin=1cm, itemsep=0.5em]
    \item[\textbf{A)}] $|p| \leq 2$ এবং $q = 1$ ($|p| \leq 2$ and $q = 1$)
    \item[\textbf{B)}] $|p| \leq 2$ এবং $q = -1$ ($|p| \leq 2$ and $q = -1$)
    \item[\textbf{C)}] $|p| > 2$ এবং $q = 1$ ($|p| > 2$ and $q = 1$)
    \item[\textbf{D)}] $|p| \leq 1$ এবং $q = 1$ ($|p| \leq 1$ and $q = 1$)
\end{itemize}

% ------------------------------------------
% ANSWER HIGHLIGHT
% ------------------------------------------
\vspace{0.5em}
\begin{tcolorbox}[answerbox]
    \textbf{\color{success} উত্তর:} A) $|p| \leq 2$ এবং $q = 1$
\end{tcolorbox}
\vspace{1em}

% ------------------------------------------
% SOLUTION SECTION
% ------------------------------------------
\begin{tcolorbox}[solutionbox={সমাধান (Solution)}]
    \noindent
    যেহেতু সমীকরণের সহগদ্বয় $p, q$ বাস্তব সংখ্যা, সেহেতু এর মূলদ্বয় অবশ্যই কাল্পনিক অনুবন্ধী যুগল হবে অথবা বাস্তব সংখ্যা হবে। 
    
    যদি মূলদ্বয় কাল্পনিক অনুবন্ধী যুগল হয়, তবে $z_2 = \bar{z}_1$। মূলদ্বয়ের গুণফল:
    \[ q = z_1 z_2 = z_1 \bar{z}_1 = |z_1|^2 \]
    দেওয়া আছে, $|z_1| = 1$। অতএব, $q = 1^2 = 1$। 
    
    আবার, মূলদ্বয়ের যোগফল:
    \[ -p = z_1 + z_2 = z_1 + \bar{z}_1 = 2 \text{Re}(z_1) \]
    আমরা জানি, যেকোনো জটিল সংখ্যার জন্য $|\text{Re}(z_1)| \leq |z_1|$। যেহেতু $|z_1| = 1$, সেহেতু $|\text{Re}(z_1)| \leq 1$। অতএব:
    \[ |-p| = 2 |\text{Re}(z_1)| \leq 2 \Rightarrow |p| \leq 2 \]
    
    যদি মূলদ্বয় বাস্তব হয়, তবে $|z_1| = |z_2| = 1$ হওয়ার কারণে মূলদ্বয় হতে পারে $1, 1$ অথবা $-1, -1$ অথবা $1, -1$। প্রথম দুই ক্ষেত্রে $q = 1$ এবং $|p| = 2 \leq 2$ যা শর্তটির সাথে সামঞ্জস্যপূর্ণ। 
    
    সুতরাং, সাধারণ ক্ষেত্রটির জন্য সঠিক উত্তরটি হলো $|p| \leq 2$ এবং $q = 1$।
\end{tcolorbox}

\vspace{1em}

% ------------------------------------------
% QUESTION SECTION
% ------------------------------------------
\begin{tcolorbox}[questionbox={প্রশ্ন (Question) 8}]
    \noindent
    আর্গন্ড সমতলে $z_1, z_2, z_3$ শীর্ষ বিশিষ্ট ত্রিভুজটি সমকোণী সমদ্বিবাহু ত্রিভুজ হওয়ার শর্ত কোনটি, যেখানে সমকোণটি $z_2$ শীর্ষবিন্দুতে অবস্থিত?
    
    \vspace{0.8em}
    \begin{english}
    \noindent\textit{\color{darkgray}
    What is the condition for a triangle with vertices $z_1, z_2, z_3$ in the Argand plane to be a right-angled isosceles triangle, with the right angle at vertex $z_2$?
    }
    \end{english}
\end{tcolorbox}

% ------------------------------------------
% OPTIONS SECTION
% ------------------------------------------
\begin{itemize}[label={}, leftmargin=1cm, itemsep=0.5em]
    \item[\textbf{A)}] $(z_1 - z_2)^2 + (z_3 - z_2)^2 = 0$
    \item[\textbf{B)}] $(z_1 - z_2)^2 = (z_3 - z_2)^2$
    \item[\textbf{C)}] $(z_1 - z_2)^2 + (z_3 - z_2)^2 = z_1 z_3$
    \item[\textbf{D)}] $(z_1 - z_3)^2 + (z_2 - z_3)^2 = 0$
\end{itemize}

% ------------------------------------------
% ANSWER HIGHLIGHT
% ------------------------------------------
\vspace{0.5em}
\begin{tcolorbox}[answerbox]
    \textbf{\color{success} উত্তর:} A) $(z_1 - z_2)^2 + (z_3 - z_2)^2 = 0$
\end{tcolorbox}
\vspace{1em}

% ------------------------------------------
% SOLUTION SECTION
% ------------------------------------------
\begin{tcolorbox}[solutionbox={সমাধান (Solution)}]
    \noindent
    যেহেতু ত্রিভুজটি সমদ্বিবাহু সমকোণী এবং সমকোণটি $z_2$ শীর্ষবিন্দুতে অবস্থিত, সেহেতু বাহু $z_2 z_1$ এবং $z_2 z_3$ এর দৈর্ঘ্য সমান এবং এদের মধ্যবর্তী কোণ $90^\circ$ বা $\pi/2$। 
    
    জটিল সমতলে ঘূর্ণন সূত্রের সাহায্যে আমরা লিখতে পারি:
    \[ (z_1 - z_2) = (z_3 - z_2) e^{\pm i\pi/2} \]
    আমরা জানি, $e^{\pm i\pi/2} = \cos(\pm\pi/2) + i\sin(\pm\pi/2) = \pm i$। অতএব:
    \[ (z_1 - z_2) = \pm i (z_3 - z_2) \]
    উভয় পক্ষকে বর্গ করে পাই:
    \[ (z_1 - z_2)^2 = (\pm i)^2 (z_3 - z_2)^2 \]
    যেহেতু $(\pm i)^2 = -1$, তাই:
    \[ (z_1 - z_2)^2 = - (z_3 - z_2)^2 \]
    \[ \Rightarrow (z_1 - z_2)^2 + (z_3 - z_2)^2 = 0 \]
\end{tcolorbox}
% ------------------------------------------
% QUESTION SECTION
% ------------------------------------------
\begin{tcolorbox}[questionbox={প্রশ্ন (Question) 18}]
    \noindent
    যদি $\omega$ এককের কাল্পনিক ঘনমূল হয়, তবে $\begin{vmatrix} x+1 & \omega & \omega^2 \\ \omega & x+\omega^2 & 1 \\ \omega^2 & 1 & x+\omega \end{vmatrix} = 0$ সমীকরণটির একটি সমাধান নিচের কোনটি?
    
    \vspace{0.8em}
    \begin{english}
    \noindent\textit{\color{darkgray}
    If $\omega$ is a non-real cube root of unity, then which of the following is a solution to the equation $\begin{vmatrix} x+1 & \omega & \omega^2 \\ \omega & x+\omega^2 & 1 \\ \omega^2 & 1 & x+\omega \end{vmatrix} = 0$?
    }
    \end{english}
\end{tcolorbox}

% ------------------------------------------
% OPTIONS SECTION
% ------------------------------------------
\begin{itemize}[label={}, leftmargin=1cm, itemsep=0.5em]
    \item[\textbf{A)}] $x = 0$
    \item[\textbf{B)}] $x = 1$
    \item[\textbf{C)}] $x = \omega$
    \item[\textbf{D)}] $x = \omega^2$
\end{itemize}

% ------------------------------------------
% ANSWER HIGHLIGHT
% ------------------------------------------
\vspace{0.5em}
\begin{tcolorbox}[answerbox]
    \textbf{\color{success} উত্তর:} A) $x = 0$
\end{tcolorbox}
\vspace{1em}

% ------------------------------------------
% SOLUTION SECTION
% ------------------------------------------
\begin{tcolorbox}[solutionbox={সমাধান (Solution)}]
    \noindent
    প্রদত্ত নির্ণায়ক সমীকরণটি হলো:
    \[ \begin{vmatrix} x+1 & \omega & \omega^2 \\ \omega & x+\omega^2 & 1 \\ \omega^2 & 1 & x+\omega \end{vmatrix} = 0 \]
    নির্ণায়কের সারিতে অপারেশন প্রয়োগ করে পাই, প্রথম সারিতে সকল সারির যোগফল প্রতিস্থাপন করি ($R_1 \rightarrow R_1 + R_2 + R_3$):
    
    \vspace{0.5em}
    নতুন প্রথম সারির উপাদানগুলো হবে:
    \begin{align*}
        R_{11} &= (x + 1) + \omega + \omega^2 = x + (1 + \omega + \omega^2) = x \quad (\text{যেহেতু } 1 + \omega + \omega^2 = 0) \\
        R_{12} &= \omega + (x + \omega^2) + 1 = x + (1 + \omega + \omega^2) = x \\
        R_{13} &= \omega^2 + 1 + (x + \omega) = x + (1 + \omega + \omega^2) = x
    \end{align*}
    
    \vspace{0.5em}
    অতএব, সমীকরণটি দাঁড়ায়:
    \[ \begin{vmatrix} x & x & x \\ \omega & x+\omega^2 & 1 \\ \omega^2 & 1 & x+\omega \end{vmatrix} = 0 \]
    প্রথম সারি হতে সাধারণ উৎপাদক $x$ বাইরে নিয়ে পাই:
    \[ x \begin{vmatrix} 1 & 1 & 1 \\ \omega & x+\omega^2 & 1 \\ \omega^2 & 1 & x+\omega \end{vmatrix} = 0 \]
    স্পষ্টতই, এই সমীকরণটির একটি নিশ্চিত সমাধান হলো $x = 0$।
\end{tcolorbox}
% ------------------------------------------
% QUESTION SECTION
% ------------------------------------------
\begin{tcolorbox}[questionbox={প্রশ্ন (Question) 10}]
    \noindent
    জটিল সংখ্যার মান নির্ণয়ের নিয়মানুযায়ী, $\sqrt{i} + \sqrt{-i}$ এর একটি সম্ভাব্য সরলীকৃত মান কোনটি?
    
    \vspace{0.8em}
    \begin{english}
    \noindent\textit{\color{darkgray}
    According to the rules of complex number evaluation, what is a possible simplified value of $\sqrt{i} + \sqrt{-i}$?
    }
    \end{english}
\end{tcolorbox}

% ------------------------------------------
% OPTIONS SECTION
% ------------------------------------------
\begin{itemize}[label={}, leftmargin=1cm, itemsep=0.5em]
    \item[\textbf{A)}] $\sqrt{2}$
    \item[\textbf{B)}] $i\sqrt{2}$
    \item[\textbf{C)}] $\pm \sqrt{2}$
    \item[\textbf{D)}] $2$
\end{itemize}

% ------------------------------------------
% ANSWER HIGHLIGHT
% ------------------------------------------
\vspace{0.5em}
\begin{tcolorbox}[answerbox]
    \textbf{\color{success} উত্তর:} C) $\pm \sqrt{2}$
\end{tcolorbox}
\vspace{1em}

% ------------------------------------------
% SOLUTION SECTION
% ------------------------------------------
\begin{tcolorbox}[solutionbox={সমাধান (Solution)}]
    \noindent
    আমরা জানি:
    \[ \sqrt{i} = \pm \frac{1}{\sqrt{2}}(1+i) \]
    \[ \sqrt{-i} = \pm \frac{1}{\sqrt{2}}(1-i) \]
    একই চিহ্নবিশিষ্ট জোড় বিবেচনা করে সমষ্টি নির্ণয় করি:
    \begin{itemize}
        \item \textbf{ধনাত্মক চিহ্ন নিয়ে পাই:} 
        \[ S = \frac{1}{\sqrt{2}}(1+i) + \frac{1}{\sqrt{2}}(1-i) = \frac{1}{\sqrt{2}}(1 + i + 1 - i) = \frac{2}{\sqrt{2}} = \sqrt{2} \]
        \item \textbf{ঋণাত্মক চিহ্ন নিয়ে পাই:} 
        \[ S = -\frac{1}{\sqrt{2}}(1+i) - \frac{1}{\sqrt{2}}(1-i) = -\frac{2}{\sqrt{2}} = -\sqrt{2} \]
    \end{itemize}
    অতএব, $\sqrt{i} + \sqrt{-i}$ এর সম্ভাব্য মান হলো $\pm \sqrt{2}$।
\end{tcolorbox}
% ------------------------------------------
% QUESTION SECTION
% ------------------------------------------
\begin{tcolorbox}[questionbox={প্রশ্ন (Question) 11}]
    \noindent
    আর্গন্ড সমতলে যেকোনো জটিল সংখ্যা $z$ এর জন্য $|z| - |z-1|$ এর সর্বোচ্চ মান কত?
    
    \vspace{0.8em}
    \begin{english}
    \noindent\textit{\color{darkgray}
    For any complex number $z$ in the Argand plane, what is the maximum value of $|z| - |z-1|$?
    }
    \end{english}
\end{tcolorbox}

% ------------------------------------------
% OPTIONS SECTION
% ------------------------------------------
\begin{itemize}[label={}, leftmargin=1cm, itemsep=0.5em]
    \item[\textbf{A)}] $0$
    \item[\textbf{B)}] $1/2$
    \item[\textbf{C)}] $1$
    \item[\textbf{D)}] $3/2$
\end{itemize}

% ------------------------------------------
% ANSWER HIGHLIGHT
% ------------------------------------------
\vspace{0.5em}
\begin{tcolorbox}[answerbox]
    \textbf{\color{success} উত্তর:} C) $1$
\end{tcolorbox}
\vspace{1em}

% ------------------------------------------
% SOLUTION SECTION
% ------------------------------------------
\begin{tcolorbox}[solutionbox={সমাধান (Solution)}]
    \noindent
    জ্যামিতিক ত্রিভুজ অসমতা (Triangle Inequality) ব্যবহার করে পাই:
    \[ |z| = |(z-1) + 1| \leq |z-1| + |1| \]
    \[ \Rightarrow |z| - |z-1| \leq 1 \]
    সমতাটি তখনই সত্য হবে যখন $z-1$ এবং $1$ একই অভিমুখে থাকবে, অর্থাৎ $z$ বাস্তব অক্ষের উপর এবং $z \geq 1$ হবে। 
    
    অতএব, $|z| - |z-1|$ এর সর্বোচ্চ মানটি হলো $1$।
\end{tcolorbox}

\vspace{1em}

% ------------------------------------------
% QUESTION SECTION
% ------------------------------------------
\begin{tcolorbox}[questionbox={প্রশ্ন (Question) 12}]
    \noindent
    যদি $z$ একক বৃত্ত $|z| = 1$ (যেখানে $z \neq 1$) এ অবস্থিত একটি জটিল সংখ্যা হয়, তবে $\text{Re}\left(\frac{1}{1-z}\right)$ এর মান কত?
    
    \vspace{0.8em}
    \begin{english}
    \noindent\textit{\color{darkgray}
    If $z$ is a complex number lying on the unit circle $|z| = 1$ (where $z \neq 1$), then what is the value of $\text{Re}\left(\frac{1}{1-z}\right)$?
    }
    \end{english}
\end{tcolorbox}

% ------------------------------------------
% OPTIONS SECTION
% ------------------------------------------
\begin{itemize}[label={}, leftmargin=1cm, itemsep=0.5em]
    \item[\textbf{A)}] $0$
    \item[\textbf{B)}] $1/2$
    \item[\textbf{C)}] $1$
    \item[\textbf{D)}] $-1/2$
\end{itemize}

% ------------------------------------------
% ANSWER HIGHLIGHT
% ------------------------------------------
\vspace{0.5em}
\begin{tcolorbox}[answerbox]
    \textbf{\color{success} উত্তর:} B) $1/2$
\end{tcolorbox}
\vspace{1em}

% ------------------------------------------
% SOLUTION SECTION
% ------------------------------------------
\begin{tcolorbox}[solutionbox={সমাধান (Solution)}]
    \noindent
    ধরি, $z = e^{i\theta} = \cos\theta + i\sin\theta$ (যেখানে $\theta \neq 2k\pi$)। প্রদত্ত রাশিটি হলো:
    \[ \frac{1}{1-z} = \frac{1}{1-\cos\theta - i\sin\theta} \]
    লব ও হরকে অনুবন্ধী দ্বারা গুণ করে পাই:
    \begin{align*}
        \frac{1}{1-z} &= \frac{1-\cos\theta + i\sin\theta}{(1-\cos\theta)^2 + \sin^2\theta} = \frac{1-\cos\theta + i\sin\theta}{1 - 2\cos\theta + \cos^2\theta + \sin^2\theta} \\
        &= \frac{1-\cos\theta + i\sin\theta}{2(1-\cos\theta)}
    \end{align*}
    বাস্তব অংশ আলাদা করে পাই:
    \[ \text{Re}\left(\frac{1}{1-z}\right) = \frac{1-\cos\theta}{2(1-\cos\theta)} = \frac{1}{2} \quad (\text{যেহেতু } \cos\theta \neq 1) \]
    অতএব, এর মান সর্বদা $1/2$।
\end{tcolorbox}

\vspace{1em}

% ------------------------------------------
% QUESTION SECTION
% ------------------------------------------
\begin{tcolorbox}[questionbox={প্রশ্ন (Question) 13}]
    \noindent
    যদি কোনো জটিল সংখ্যা $z$ এর জন্য $|z - i \text{Re}(z)| = |z - \text{Im}(z)|$ সমীকরণটি সত্য হয়, তবে কার্তেসীয় সমতলে সঞ্চারপথটি কেমন হবে?
    
    \vspace{0.8em}
    \begin{english}
    \noindent\textit{\color{darkgray}
    If for a complex number $z$, the equation $|z - i \text{Re}(z)| = |z - \text{Im}(z)|$ is true, then what is the locus in the cartesian plane?
    }
    \end{english}
\end{tcolorbox}

% ------------------------------------------
% OPTIONS SECTION
% ------------------------------------------
\begin{itemize}[label={}, leftmargin=1cm, itemsep=0.5em]
    \item[\textbf{A)}] $y = \pm x$
    \item[\textbf{B)}] $y = x$
    \item[\textbf{C)}] $y = 0$
    \item[\textbf{D)}] $x = 0$
\end{itemize}

% ------------------------------------------
% ANSWER HIGHLIGHT
% ------------------------------------------
\vspace{0.5em}
\begin{tcolorbox}[answerbox]
    \textbf{\color{success} উত্তর:} A) $y = \pm x$
\end{tcolorbox}
\vspace{1em}

% ------------------------------------------
% SOLUTION SECTION
% ------------------------------------------
\begin{tcolorbox}[solutionbox={সমাধান (Solution)}]
    \noindent
    ধরি, $z = x + iy$। তাহলে, $\text{Re}(z) = x$ এবং $\text{Im}(z) = y$।
    
    বামপক্ষ: 
    \[ |z - i\text{Re}(z)| = |x + iy - ix| = |x + i(y-x)| = \sqrt{x^2 + (y-x)^2} \]
    ডানপক্ষ: 
    \[ |z - \text{Im}(z)| = |x + iy - y| = |(x-y) + iy| = \sqrt{(x-y)^2 + y^2} \]
    শর্তানুসারে এদের বর্গ করে সমীকৃত করি:
    \begin{align*}
        x^2 + (y-x)^2 &= (x-y)^2 + y^2 \\
        \Rightarrow x^2 + y^2 - 2xy + x^2 &= x^2 - 2xy + y^2 + y^2 \\
        \Rightarrow 2x^2 - 2xy + y^2 &= x^2 - 2xy + 2y^2 \\
        \Rightarrow x^2 &= y^2 \Rightarrow y = \pm x
    \end{align*}
\end{tcolorbox}

\vspace{1em}

% ------------------------------------------
% QUESTION SECTION
% ------------------------------------------
\begin{tcolorbox}[questionbox={প্রশ্ন (Question) 14}]
    \noindent
    যদি জটিল সংখ্যা $z$ এর পরমমান $|z| = 1$ এবং কাল্পনিক অংশ $\text{Im}(z) > 0$ হয়, তবে $\arg\left(\frac{z-1}{z+1}\right)$ এর মান কত?
    
    \vspace{0.8em}
    \begin{english}
    \noindent\textit{\color{darkgray}
    If for a complex number $z$, the modulus is $|z| = 1$ and the imaginary part is $\text{Im}(z) > 0$, then what is the value of $\arg\left(\frac{z-1}{z+1}\right)$?
    }
    \end{english}
\end{tcolorbox}

% ------------------------------------------
% OPTIONS SECTION
% ------------------------------------------
\begin{itemize}[label={}, leftmargin=1cm, itemsep=0.5em]
    \item[\textbf{A)}] $0$
    \item[\textbf{B)}] $\pi$
    \item[\textbf{C)}] $\pi/2$
    \item[\textbf{D)}] $-\pi/2$
\end{itemize}

% ------------------------------------------
% ANSWER HIGHLIGHT
% ------------------------------------------
\vspace{0.5em}
\begin{tcolorbox}[answerbox]
    \textbf{\color{success} উত্তর:} C) $\pi/2$
\end{tcolorbox}
\vspace{1em}

% ------------------------------------------
% SOLUTION SECTION
% ------------------------------------------
\begin{tcolorbox}[solutionbox={সমাধান (Solution)}]
    \noindent
    যেহেতু $|z| = 1$ এবং $\text{Im}(z) > 0$, তাই আমরা অয়লার আকারে লিখতে পারি: $z = e^{i\theta}$, যেখানে $\theta \in (0, \pi)$। 
    
    এখন, অনুপাতটি সরল করি:
    \begin{align*}
        \frac{z-1}{z+1} &= \frac{e^{i\theta}-1}{e^{i\theta}+1} = \frac{\cos\theta - 1 + i\sin\theta}{\cos\theta + 1 + i\sin\theta} \\
        &= \frac{-2\sin^2(\theta/2) + 2i\sin(\theta/2)\cos(\theta/2)}{2\cos^2(\theta/2) + 2i\sin(\theta/2)\cos(\theta/2)} \\
        &= \frac{2i\sin(\theta/2)[\cos(\theta/2) + i\sin(\theta/2)]}{2\cos(\theta/2)[\cos(\theta/2) + i\sin(\theta/2)]} = i\tan\left(\frac{\theta}{2}\right)
    \end{align*}
    যেহেতু $\theta \in (0, \pi)$, সেহেতু $\frac{\theta}{2} \in (0, \pi/2)$ এবং $\tan(\theta/2) > 0$। অতএব, $i\tan(\theta/2)$ একটি বিশুদ্ধ ধনাত্মক কাল্পনিক সংখ্যা। 
    
    আমরা জানি, যেকোনো ধনাত্মক কাল্পনিক সংখ্যার আর্গুমেন্ট সর্বদা $\pi/2$। সুতরাং, $\arg\left(\frac{z-1}{z+1}\right) = \frac{\pi}{2}$।
\end{tcolorbox}

\vspace{1em}

% ------------------------------------------
% QUESTION SECTION
% ------------------------------------------
\begin{tcolorbox}[questionbox={প্রশ্ন (Question) 15}]
    \noindent
    যদি $\omega$ এককের একটি কাল্পনিক ঘনমূল হয়, তবে $\arg(i\omega) + \arg(i\omega^2)$ এর সরলীকৃত মান কত?
    
    \vspace{0.8em}
    \begin{english}
    \noindent\textit{\color{darkgray}
    If $\omega$ is a non-real cube root of unity, then what is the simplified value of $\arg(i\omega) + \arg(i\omega^2)$?
    }
    \end{english}
\end{tcolorbox}

% ------------------------------------------
% OPTIONS SECTION
% ------------------------------------------
\begin{itemize}[label={}, leftmargin=1cm, itemsep=0.5em]
    \item[\textbf{A)}] $0$
    \item[\textbf{B)}] $\pi/2$
    \item[\textbf{C)}] $\pi$
    \item[\textbf{D)}] $-\pi/2$
\end{itemize}

% ------------------------------------------
% ANSWER HIGHLIGHT
% ------------------------------------------
\vspace{0.5em}
\begin{tcolorbox}[answerbox]
    \textbf{\color{success} উত্তর:} C) $\pi$
\end{tcolorbox}
\vspace{1em}

% ------------------------------------------
% SOLUTION SECTION
% ------------------------------------------
\begin{tcolorbox}[solutionbox={সমাধান (Solution)}]
    \noindent
    আমরা জানি, এককের কাল্পনিক ঘনমূলের মুখ্য আর্গুমেন্ট যথাক্রমে: $\arg(\omega) = \frac{2\pi}{3}$ এবং $\arg(\omega^2) = -\frac{2\pi}{3}$। 
    
    এখন আর্গুমেন্টের ধর্মাবলি প্রয়োগ করে পাই:
    \[ \arg(i\omega) = \arg(i) + \arg(\omega) = \frac{\pi}{2} + \frac{2\pi}{3} = \frac{7\pi}{6} \]
    যেহেতু $\frac{7\pi}{6} > \pi$, তাই এর মুখ্য মান হবে: $\frac{7\pi}{6} - 2\pi = -\frac{5\pi}{6}$। 
    
    আবার, 
    \[ \arg(i\omega^2) = \arg(i) + \arg(\omega^2) = \frac{\pi}{2} - \frac{2\pi}{3} = -\frac{\pi}{6} \quad (\text{যা } (-\pi, \pi] \text{ সীমার মধ্যে রয়েছে}) \]
    অতএব, এদের সমষ্টি:
    \[ \arg(i\omega) + \arg(i\omega^2) = -\frac{5\pi}{6} - \frac{\pi}{6} = -\pi \]
    যেহেতু আর্গন্ড সমতলে $-\pi$ এবং $\pi$ একই অভিমুখে নির্দেশ করে এবং মুখ্য মানের সীমা অনুসারে ধনাত্মক দিকটি গৃহীত হয়: অতএব, নির্ণেয় মান $\pi$।
\end{tcolorbox}
% ------------------------------------------
% QUESTION SECTION
% ------------------------------------------
\begin{tcolorbox}[questionbox={প্রশ্ন (Question) 16}]
    \noindent
    যদি জটিল সংখ্যা $z$ অসমতা $|z - 1| + |z + 3| \leq 8$ সিদ্ধ করে, তবে $|z - 4|$ এর মানের পরম পাল্লা (range) নিচের কোনটি?
    
    \vspace{0.8em}
    \begin{english}
    \noindent\textit{\color{darkgray}
    If the complex number $z$ satisfies the inequality $|z - 1| + |z + 3| \leq 8$, then the range of values of $|z - 4|$ is:
    }
    \end{english}
\end{tcolorbox}

% ------------------------------------------
% OPTIONS SECTION
% ------------------------------------------
\begin{itemize}[label={}, leftmargin=1cm, itemsep=0.5em]
    \item[\textbf{A)}] 12
    \item[\textbf{B)}] 13
    \item[\textbf{C)}] 34
    \item[\textbf{D)}] 3
\end{itemize}

% ------------------------------------------
% ANSWER HIGHLIGHT
% ------------------------------------------
\vspace{0.5em}
\begin{tcolorbox}[answerbox]
    \textbf{\color{success} উত্তর:} B) 13
\end{tcolorbox}
\vspace{1em}

% ------------------------------------------
% SOLUTION SECTION
% ------------------------------------------
\begin{tcolorbox}[solutionbox={সমাধান (Solution)}]
    \noindent
    অসমতা $|z - 1| + |z + 3| \leq 8$ আর্গন্ড চিত্রে একটি উপবৃত্তের সীমানা ও অভ্যন্তরীণ অঞ্চল প্রকাশ করে, যার উপকেন্দ্রদ্বয় (foci) $S_1 = 1 \equiv (1, 0)$ এবং $S_2 = -3 \equiv (-3, 0)$ বিন্দুতে অবস্থিত। 
    
    উপবৃত্তটির কেন্দ্র $C = \frac{1 + (-3)}{2} = -1 \equiv (-1, 0)$। উপকেন্দ্রদ্বয়ের মধ্যবর্তী দূরত্ব $2ae = 1 - (-3) = 4$। বৃহৎ অক্ষের দৈর্ঘ্য $2a = 8 \Rightarrow a = 4$ (বৃহৎ অর্ধ-অক্ষ)। যেহেতু $a = 4$ এবং কেন্দ্র $(-1, 0)$, তাই উপবৃত্তের প্রধান অক্ষটি বাস্তব অক্ষ বরাবর বিস্তৃত যার শীর্ষবিন্দুদ্বয় যথাক্রমে $(-1-4, 0) = (-5, 0)$ এবং $(-1+4, 0) = (3, 0)$। 
    
    আমরা $|z - 4|$ এর পাল্লা নির্ণয় করতে চাই, যা মূলত উপবৃত্তের যেকোনো বিন্দু $z$ হতে স্থির বিন্দু $P(4, 0)$ এর দূরত্ব। যেহেতু $P(4, 0)$ বিন্দুটি উপবৃত্তের বাইরে বাস্তব অক্ষের উপর ডানদিকে অবস্থিত:
    \begin{enumerate}
        \item সর্বনিম্ন দূরত্ব হবে $P(4, 0)$ হতে উপবৃত্তের সবচেয়ে কাছের শীর্ষবিন্দু $(3, 0)$ এর দূরত্ব: $d_{\text{min}} = |4 - 3| = 1$।
        \item সর্বোচ্চ দূরত্ব হবে $P(4, 0)$ হতে উপবৃত্তের সবচেয়ে দূরের শীর্ষবিন্দু $(-5, 0)$ এর দূরত্ব: $d_{\text{max}} = |4 - (-5)| = 9$।
    \end{enumerate}
    অতএব, $|z - 4|$ এর পরম পাল্লা হলো 13।
\end{tcolorbox}

\vspace{1em}

% ------------------------------------------
% QUESTION SECTION
% ------------------------------------------
\begin{tcolorbox}[questionbox={প্রশ্ন (Question) 17}]
    \noindent
    যদি $\left(\frac{3}{2} + \frac{\sqrt{3}}{2}i\right)^{50} = 3^{25}(x + iy)$ হয় (যেখানে $x, y \in \mathbb{R}$), তবে ক্রোমজোড় $(x, y)$ এর মান নিচের কোনটি?
    
    \vspace{0.8em}
    \begin{english}
    \noindent\textit{\color{darkgray}
    If $\left(\frac{3}{2} + \frac{\sqrt{3}}{2}i\right)^{50} = 3^{25}(x + iy)$ (where $x, y \in \mathbb{R}$), then the ordered pair $(x, y)$ is:
    }
    \end{english}
\end{tcolorbox}

% ------------------------------------------
% OPTIONS SECTION
% ------------------------------------------
\begin{itemize}[label={}, leftmargin=1cm, itemsep=0.5em]
    \item[\textbf{A)}] $\left(\frac{1}{2}, \frac{\sqrt{3}}{2}\right)$
    \item[\textbf{B)}] $\left(\frac{\sqrt{3}}{2}, \frac{1}{2}\right)$
    \item[\textbf{C)}] $\left(-\frac{1}{2}, \frac{\sqrt{3}}{2}\right)$
    \item[\textbf{D)}] $\left(\frac{1}{2}, -\frac{\sqrt{3}}{2}\right)$
\end{itemize}

% ------------------------------------------
% ANSWER HIGHLIGHT
% ------------------------------------------
\vspace{0.5em}
\begin{tcolorbox}[answerbox]
    \textbf{\color{success} উত্তর:} A) $\left(\frac{1}{2}, \frac{\sqrt{3}}{2}\right)$
\end{tcolorbox}
\vspace{1em}

% ------------------------------------------
% SOLUTION SECTION
% ------------------------------------------
\begin{tcolorbox}[solutionbox={সমাধান (Solution)}]
    \noindent
    ধরি, বন্ধনীভুক্ত জটিল সংখ্যাটি $z = \frac{3}{2} + \frac{\sqrt{3}}{2}i$। একে পোলার বা অয়লার আকারে রূপান্তর করি:
    \[ z = \sqrt{3}\left(\frac{\sqrt{3}}{2} + \frac{1}{2}i\right) = \sqrt{3} e^{i\pi/6} \]
    এখন ৫০-তম ঘাত প্রয়োগ করে পাই:
    \[ z^{50} = (\sqrt{3} e^{i\pi/6})^{50} = (\sqrt{3})^{50} e^{i 50\pi/6} = 3^{25} e^{i 25\pi/3} \]
    আর্গুমেন্টটিকে কোণীয় ব্যবধানে সরল করে পাই:
    \[ \frac{25\pi}{3} = 8\pi + \frac{\pi}{3} \equiv \frac{\pi}{3} \pmod{2\pi} \]
    অতএব:
    \[ e^{i 25\pi/3} = e^{i\pi/3} = \cos\left(\frac{\pi}{3}\right) + i\sin\left(\frac{\pi}{3}\right) = \frac{1}{2} + i\frac{\sqrt{3}}{2} \]
    সুতরাং:
    \[ z^{50} = 3^{25}\left(\frac{1}{2} + i\frac{\sqrt{3}}{2}\right) \]
    প্রদত্ত সমীকরণের সাথে তুলনা করে পাই: $x = \frac{1}{2}$ এবং বা $y = \frac{\sqrt{3}}{2}$। অতএব, ক্রোমজোড় $(x, y) = \left(\frac{1}{2}, \frac{\sqrt{3}}{2}\right)$।
\end{tcolorbox}

\vspace{1em}

% ------------------------------------------
% QUESTION SECTION
% ------------------------------------------
\begin{tcolorbox}[questionbox={প্রশ্ন (Question) 18}]
    \noindent
    যদি $z_1$ এবং $z_2$ দুটি জটিল সংখ্যা হয় যা $\left|\frac{z_1 - 3z_2}{3 - z_1 \bar{z}_2}\right| = 1$ সমীকরণটি সিদ্ধ করে এবং $|z_2| \neq 1$ হয়, তবে $|z_1|$ এর সঠিক মান কত?
    
    \vspace{0.8em}
    \begin{english}
    \noindent\textit{\color{darkgray}
    If $z_1$ and $z_2$ are two complex numbers satisfying the equation $\left|\frac{z_1 - 3z_2}{3 - z_1 \bar{z}_2}\right| = 1$ and $|z_2| \neq 1$, then what is the correct value of $|z_1|$?
    }
    \end{english}
\end{tcolorbox}

% ------------------------------------------
% OPTIONS SECTION
% ------------------------------------------
\begin{itemize}[label={}, leftmargin=1cm, itemsep=0.5em]
    \item[\textbf{A)}] $1$
    \item[\textbf{B)}] $2$
    \item[\textbf{C)}] $3$
    \item[\textbf{D)}] $9$
\end{itemize}

% ------------------------------------------
% ANSWER HIGHLIGHT
% ------------------------------------------
\vspace{0.5em}
\begin{tcolorbox}[answerbox]
    \textbf{\color{success} উত্তর:} C) $3$
\end{tcolorbox}
\vspace{1em}

% ------------------------------------------
% SOLUTION SECTION
% ------------------------------------------
\begin{tcolorbox}[solutionbox={সমাধান (Solution)}]
    \noindent
    প্রদত্ত সমীকরণটি থেকে পাই:
    \[ |z_1 - 3z_2| = |3 - z_1 \bar{z}_2| \]
    উভয় পক্ষকে বর্গ করে পাই:
    \[ |z_1 - 3z_2|^2 = |3 - z_1 \bar{z}_2|^2 \]
    জটিল সংখ্যার বর্গ গুণের ধর্মাবলি প্রয়োগ করে পাই ($|Z|^2 = Z\bar{Z}$):
    \begin{align*}
        (z_1 - 3z_2)(\bar{z}_1 - 3\bar{z}_2) &= (3 - z_1 \bar{z}_2)(3 - \bar{z}_1 z_2) \\
        \Rightarrow z_1 \bar{z}_1 - 3z_1 \bar{z}_2 - 3\bar{z}_1 z_2 + 9z_2 \bar{z}_2 &= 9 - 3z_1 \bar{z}_2 - 3\bar{z}_1 z_2 + z_1 \bar{z}_1 z_2 \bar{z}_2
    \end{align*}
    উভয় পক্ষ হতে একই কাল্পনিক পদ $-3(z_1 \bar{z}_2 + \bar{z}_1 z_2)$ বর্জন করে পাই:
    \[ |z_1|^2 + 9|z_2|^2 = 9 + |z_1|^2 |z_2|^2 \]
    সবগুলো পদ একপাশে নিয়ে উৎপাদকে বিশ্লেষণ করি:
    \begin{align*}
        |z_1|^2 - |z_1|^2 |z_2|^2 - 9 + 9|z_2|^2 &= 0 \\
        \Rightarrow |z_1|^2(1 - |z_2|^2) - 9(1 - |z_2|^2) &= 0 \\
        \Rightarrow (|z_1|^2 - 9)(1 - |z_2|^2) &= 0
    \end{align*}
    যেহেতু দেওয়া আছে $|z_2| \neq 1 \Rightarrow 1 - |z_2|^2 \neq 0$। অতএব:
    \[ |z_1|^2 - 9 = 0 \Rightarrow |z_1|^2 = 9 \Rightarrow |z_1| = 3 \]
\end{tcolorbox}

\vspace{1em}

% ------------------------------------------
% QUESTION SECTION
% ------------------------------------------
\begin{tcolorbox}[questionbox={প্রশ্ন (Question) 19}]
    \noindent
    জটিল সমীকরণ $z^5 = \bar{z}$ এবং $\arg(z + 1) = \frac{\pi}{6}$ উভয় শর্তকে একসাথে সিদ্ধ করে এমন একক সমাধানবিশিষ্ট জটিল সংখ্যা $z$ নিচের কোনটি?
    
    \vspace{0.8em}
    \begin{english}
    \noindent\textit{\color{darkgray}
    What is the unique complex number $z$ that simultaneously satisfies the equations $z^5 = \bar{z}$ and $\arg(z + 1) = \frac{\pi}{6}$?
    }
    \end{english}
\end{tcolorbox}

% ------------------------------------------
% OPTIONS SECTION
% ------------------------------------------
\begin{itemize}[label={}, leftmargin=1cm, itemsep=0.5em]
    \item[\textbf{A)}] $\frac{1}{2} + i\frac{\sqrt{3}}{2}$
    \item[\textbf{B)}] $\frac{\sqrt{3}}{2} + \frac{1}{2}i$
    \item[\textbf{C)}] $-\frac{1}{2} + i\frac{\sqrt{3}}{2}$
    \item[\textbf{D)}] $\frac{1}{2} - i\frac{\sqrt{3}}{2}$
\end{itemize}

% ------------------------------------------
% ANSWER HIGHLIGHT
% ------------------------------------------
\vspace{0.5em}
\begin{tcolorbox}[answerbox]
    \textbf{\color{success} উত্তর:} A) $\frac{1}{2} + i\frac{\sqrt{3}}{2}$
\end{tcolorbox}
\vspace{1em}

% ------------------------------------------
% SOLUTION SECTION
% ------------------------------------------
\begin{tcolorbox}[solutionbox={সমাধান (Solution)}]
    \noindent
    প্রথম সমীকরণ থেকে পাই: $z^5 = \bar{z}$
    
    উভয় পক্ষে মডুলাস নিয়ে পাই: $|z|^5 = |\bar{z}| = |z|$। যেহেতু কোণটির আর্গুমেন্ট সংজ্ঞায়িত, তাই বা $z \neq 0$। অতএব, $|z|^4 = 1 \Rightarrow |z| = 1$। 
    
    এখন প্রথম সমীকরণের উভয় পক্ষকে $z$ দ্বারা গুণ করে পাই: $z^6 = z \bar{z} = |z|^2 = 1$
    
    সুতরাং, বা $z$ হলো এককের ৬-তম মূলসমূহ। অর্থাৎ, $z = e^{i k\pi/3}$ (যেখানে $k = 0, 1, 2, 3, 4, 5$)। 
    
    দ্বিতীয় সমীকরণ $\arg(z + 1) = \frac{\pi}{6}$ দ্বারা একটি রশ্মি নির্দেশিত হয় যা $(-1, 0)$ বিন্দু হতে শুরু হয়ে বাস্তব অক্ষের ধনাত্মক দিকের সাথে $30^\circ$ কোণ উৎপন্ন করে। ধরি, বা রশ্মি বরাবর সমীকরণ:
    \[ z + 1 = r e^{i\pi/6} = r\left(\frac{\sqrt{3}}{2} + \frac{1}{2}i\right) \quad (\text{যেখানে } r > 0) \]
    \[ \Rightarrow z = \left(\frac{r\sqrt{3}}{2} - 1\right) + i\frac{r}{2} \]
    যেহেতু $|z| = 1$, তাই:
    \begin{align*}
        \left(\frac{r\sqrt{3}}{2} - 1\right)^2 + \left(\frac{r}{2}\right)^2 &= 1 \\
        \Rightarrow \frac{3r^2}{4} - r\sqrt{3} + 1 + \frac{r^2}{4} &= 1 \\
        \Rightarrow r^2 - r\sqrt{3} &= 0
    \end{align*}
    যেহেতু $r > 0$, সেহেতু $r = \sqrt{3}$। তাহলে বা $z$ এর মান হবে:
    \[ z = \left(\frac{3}{2} - 1\right) + i\frac{\sqrt{3}}{2} = \frac{1}{2} + i\frac{\sqrt{3}}{2} = e^{i\pi/3} \]
    এটি $z^6 = 1$ এর $k=1$ এর সাথে হুবহু মিলে যায়। অতএব, নির্ণেয় জটিল সংখ্যাটি হলো $\frac{1}{2} + i\frac{\sqrt{3}}{2}$।
\end{tcolorbox}

\vspace{1em}

% ------------------------------------------
% QUESTION SECTION
% ------------------------------------------
\begin{tcolorbox}[questionbox={প্রশ্ন (Question) 20}]
    \noindent
    যদি $z = \frac{1}{2}(\sqrt{3} - i)$ এবং $(z^{101} + i^{101})^{124} = z^n$ সমীকরণটি সিদ্ধকারী ক্ষুদ্রতম ধনাত্মক পূর্ণসাংখ্যিক মান $n = K$ হয়, তবে বা $\frac{24}{K}$ এর সঠিক মান কত?
    
    \vspace{0.8em}
    \begin{english}
    \noindent\textit{\color{darkgray}
    If $z = \frac{1}{2}(\sqrt{3} - i)$ and the least positive integral value of $n$ satisfying the equation $(z^{101} + i^{101})^{124} = z^n$ is $n = K$, then what is the correct value of $\frac{24}{K}$?
    }
    \end{english}
\end{tcolorbox}

% ------------------------------------------
% OPTIONS SECTION
% ------------------------------------------
\begin{itemize}[label={}, leftmargin=1cm, itemsep=0.5em]
    \item[\textbf{A)}] $3$
    \item[\textbf{B)}] $4$
    \item[\textbf{C)}] $6$
    \item[\textbf{D)}] $12$
\end{itemize}

% ------------------------------------------
% ANSWER HIGHLIGHT
% ------------------------------------------
\vspace{0.5em}
\begin{tcolorbox}[answerbox]
    \textbf{\color{success} উত্তর:} C) $6$
\end{tcolorbox}
\vspace{1em}

% ------------------------------------------
% SOLUTION SECTION
% ------------------------------------------
\begin{tcolorbox}[solutionbox={সমাধান (Solution)}]
    \noindent
    আমরা জানি, $z = \frac{\sqrt{3}}{2} - \frac{1}{2}i = e^{-i\pi/6}$। তাহলে:
    \[ z^{101} = e^{-i 101\pi/6} \]
    কোণটিকে পর্যায়বৃত্ততা অনুসারে ছোট করে পাই:
    \[ -\frac{101\pi}{6} = -16\pi - \frac{5\pi}{6} \equiv -\frac{5\pi}{6} \pmod{2\pi} \]
    \[ \Rightarrow z^{101} = \cos\left(-\frac{5\pi}{6}\right) + i\sin\left(-\frac{5\pi}{6}\right) = -\frac{\sqrt{3}}{2} - \frac{1}{2}i \]
    আমরা জানি, $i^{101} = i$। সুতরাং:
    \[ z^{101} + i^{101} = -\frac{\sqrt{3}}{2} - \frac{1}{2}i + i = -\frac{\sqrt{3}}{2} + \frac{1}{2}i = e^{i 5\pi/6} \]
    এখন ঘাত ১২৪ প্রয়োগ করে পাই:
    \[ (e^{i 5\pi/6})^{124} = e^{i 620\pi/6} = e^{i 310\pi/3} \]
    কোণটিকে আরও সরল করি:
    \[ \frac{310\pi}{3} = 102\pi + \frac{4\pi}{3} \equiv \frac{4\pi}{3} \equiv -\frac{2\pi}{3} \pmod{2\pi} \]
    তাহলে সমীকরণটি দাঁড়ায়:
    \begin{align*}
        z^n &= (e^{-i\pi/6})^n = e^{-i n\pi/6} = e^{-i 2\pi/3} \\
        &\Rightarrow \frac{n\pi}{6} = \frac{2\pi}{3} \Rightarrow n = 4
    \end{align*}
    অতএব, ক্ষুদ্রতম ধনাত্মক পূর্ণসাংখ্যিক মান $K = 4$। সুতরাং, $\frac{24}{K} = \frac{24}{4} = 6$।
\end{tcolorbox}

\vspace{1em}

% ------------------------------------------
% QUESTION SECTION
% ------------------------------------------
\begin{tcolorbox}[questionbox={প্রশ্ন (Question) 21}]
    \noindent
    যদি সমীকরণ $(\sqrt{3} + i)^{100} = 2^{99}(p + iq)$ সত্য হয় (যেখানে $p, q \in \mathbb{R}$), তবে বা $p$ এবং $q$ মূল বিশিষ্ট দ্বিঘাত সমীকরণটি নিচের কোনটি?
    
    \vspace{0.8em}
    \begin{english}
    \noindent\textit{\color{darkgray}
    If the equation $(\sqrt{3} + i)^{100} = 2^{99}(p + iq)$ is true (where $p, q \in \mathbb{R}$), then the quadratic equation whose roots are $p$ and $q$ is:
    }
    \end{english}
\end{tcolorbox}

% ------------------------------------------
% OPTIONS SECTION
% ------------------------------------------
\begin{itemize}[label={}, leftmargin=1cm, itemsep=0.5em]
    \item[\textbf{A)}] বা $x^2 - (\sqrt{3}-1)x - \sqrt{3} = 0$
    \item[\textbf{B)}] বা $x^2 + (\sqrt{3}-1)x - \sqrt{3} = 0$
    \item[\textbf{C)}] বা $x^2 - (\sqrt{3}+1)x + \sqrt{3} = 0$
    \item[\textbf{D)}] বা $x^2 + (\sqrt{3}+1)x + \sqrt{3} = 0$
\end{itemize}

% ------------------------------------------
% ANSWER HIGHLIGHT
% ------------------------------------------
\vspace{0.5em}
\begin{tcolorbox}[answerbox]
    \textbf{\color{success} উত্তর:} A) $x^2 - (\sqrt{3}-1)x - \sqrt{3} = 0$
\end{tcolorbox}
\vspace{1em}

% ------------------------------------------
% SOLUTION SECTION
% ------------------------------------------
\begin{tcolorbox}[solutionbox={সমাধান (Solution)}]
    \noindent
    ধরি, জটিল সংখ্যা $Z = \sqrt{3} + i$। একে অয়লারের আকারে লিখি:
    \[ Z = 2\left(\frac{\sqrt{3}}{2} + \frac{1}{2}i\right) = 2 e^{i\pi/6} \]
    ১০০-তম ঘাত প্রয়োগ করি:
    \[ Z^{100} = (2 e^{i\pi/6})^{100} = 2^{100} e^{i 100\pi/6} = 2^{100} e^{i 50\pi/3} \]
    কোণটিকে সরলীকরণ করে পাই:
    \[ \frac{50\pi}{3} = 16\pi + \frac{2\pi}{3} \equiv \frac{2\pi}{3} \pmod{2\pi} \]
    তাহলে:
    \[ e^{i 50\pi/3} = e^{i 2\pi/3} = \cos\left(\frac{2\pi}{3}\right) + i\sin\left(\frac{2\pi}{3}\right) = -\frac{1}{2} + i\frac{\sqrt{3}}{2} \]
    এখন সমীকরণে বসাই:
    \[ Z^{100} = 2^{100}\left(-\frac{1}{2} + i\frac{\sqrt{3}}{2}\right) = 2^{99}(-1 + i\sqrt{3}) \]
    প্রদত্ত রাশির সাথে তুলনা করি:
    \[ 2^{99}(p + iq) = 2^{99}(-1 + i\sqrt{3}) \]
    সহগ সমীকৃত করে পাই: $p = -1$ এবং $q = \sqrt{3}$। 
    
    এখন, বা $p$ এবং $q$ মূলদ্বয় বিশিষ্ট দ্বিঘাত সমীকরণটি হবে:
    \begin{align*}
        x^2 - (p + q)x + pq &= 0 \\
        \Rightarrow x^2 - (\sqrt{3} - 1)x + (-1 \times \sqrt{3}) &= 0 \\
        \Rightarrow x^2 - (\sqrt{3}-1)x - \sqrt{3} &= 0
    \end{align*}
\end{tcolorbox}
% ------------------------------------------
% QUESTION SECTION
% ------------------------------------------
\begin{tcolorbox}[questionbox={প্রশ্ন (Question) 22}]
    \noindent
    যদি বা $a, x, y$ বাস্তব সংখ্যা হয় (যেখানে $a \neq 0$), $z = x + iy$ এবং $\frac{z-i}{z+1} = ia$ সমীকরণটি সিদ্ধ করে, তবে $\left(x + \frac{1}{2}\right)^2 + \left(y - \frac{1}{2}\right)^2$ এর মান কত?
    
    \vspace{0.8em}
    \begin{english}
    \noindent\textit{\color{darkgray}
    If $a, x, y$ are real numbers (with $a \neq 0$), $z = x + iy$ and $\frac{z-i}{z+1} = ia$ is satisfied, then the value of $\left(x + \frac{1}{2}\right)^2 + \left(y - \frac{1}{2}\right)^2$ is:
    }
    \end{english}
\end{tcolorbox}

% ------------------------------------------
% OPTIONS SECTION
% ------------------------------------------
\begin{itemize}[label={}, leftmargin=1cm, itemsep=0.5em]
    \item[\textbf{A)}] বা $1/4$
    \item[\textbf{B)}] বা $1/2$
    \item[\textbf{C)}] বা $1$
    \item[\textbf{D)}] বা $2$
\end{itemize}

% ------------------------------------------
% ANSWER HIGHLIGHT
% ------------------------------------------
\vspace{0.5em}
\begin{tcolorbox}[answerbox]
    \textbf{\color{success} উত্তর:} B) $1/2$
\end{tcolorbox}
\vspace{1em}

% ------------------------------------------
% SOLUTION SECTION
% ------------------------------------------
\begin{tcolorbox}[solutionbox={সমাধান (Solution)}]
    \noindent
    দেওয়া আছে, $\frac{z-i}{z+1} = ia$। যেহেতু $a$ বাস্তব সংখ্যা, সেহেতু $ia$ একটি বিশুদ্ধ কাল্পনিক সংখ্যা। অতএব, $\frac{z-i}{z+1}$ এর বাস্তব অংশ অবশ্যই শূন্য হতে হবে। 
    
    ধরি, বা $z = x+iy$। 
    \[ \frac{z-i}{z+1} = \frac{x+i(y-1)}{(x+1)+iy} = \frac{[x+i(y-1)][(x+1)-iy]}{(x+1)^2+y^2} \]
    লবের বাস্তব অংশ বের করার জন্য স্বাভাবিক বীজগণিতীয় গুণ করি:
    \[ \text{Re}(\text{Lnumerator}) = x(x+1) + y(y-1) \]
    যেহেতু সম্পূর্ণ রাশির বাস্তব অংশ শূন্য, তাই:
    \begin{align*}
        x(x+1) + y(y-1) &= 0 \\
        \Rightarrow x^2 + x + y^2 - y &= 0
    \end{align*}
    এবার বর্গ পূর্ণ করার জন্য সাজিয়ে পাই:
    \begin{align*}
        \left(x^2 + x + \frac{1}{4}\right) + \left(y^2 - y + \frac{1}{4}\right) &= \frac{1}{4} + \frac{1}{4} \\
        \Rightarrow \left(x + \frac{1}{2}\right)^2 + \left(y - \frac{1}{2}\right)^2 &= \frac{1}{2}
    \end{align*}
\end{tcolorbox}

\vspace{1em}

% ------------------------------------------
% QUESTION SECTION
% ------------------------------------------
\begin{tcolorbox}[questionbox={প্রশ্ন (Question) 23}]
    \noindent
    জটিল সংখ্যার ধর্মাবলি বিশ্লেষণ করে $4 + 5\left(-\frac{1}{2} + \frac{i\sqrt{3}}{2}\right)^{334} + 3\left(-\frac{1}{2} + \frac{i\sqrt{3}}{2}\right)^{365}$ রাশিটির সঠিক মান নির্ণয় করো।
    
    \vspace{0.8em}
    \begin{english}
    \noindent\textit{\color{darkgray}
    Analyzing the properties of complex numbers, determine the correct value of the expression $4 + 5\left(-\frac{1}{2} + \frac{i\sqrt{3}}{2}\right)^{334} + 3\left(-\frac{1}{2} + \frac{i\sqrt{3}}{2}\right)^{365}$.
    }
    \end{english}
\end{tcolorbox}

% ------------------------------------------
% OPTIONS SECTION
% ------------------------------------------
\begin{itemize}[label={}, leftmargin=1cm, itemsep=0.5em]
    \item[\textbf{A)}] বা $i\sqrt{3}$
    \item[\textbf{B)}] বা $-i\sqrt{3}$
    \item[\textbf{C)}] বা $\sqrt{3}$
    \item[\textbf{D)}] বা $-\sqrt{3}$
\end{itemize}

% ------------------------------------------
% ANSWER HIGHLIGHT
% ------------------------------------------
\vspace{0.5em}
\begin{tcolorbox}[answerbox]
    \textbf{\color{success} উত্তর:} A) $i\sqrt{3}$
\end{tcolorbox}
\vspace{1em}

% ------------------------------------------
% SOLUTION SECTION
% ------------------------------------------
\begin{tcolorbox}[solutionbox={সমাধান (Solution)}]
    \noindent
    আমরা জানি, এককের কাল্পনিক ঘনমূল $\omega = -\frac{1}{2} + i\frac{\sqrt{3}}{2}$। তাহলে প্রদত্ত রাশিটি দাঁড়ায়:
    \[ P = 4 + 5\omega^{334} + 3\omega^{365} \]
    আমরা জানি, $\omega^3 = 1$। ৩ দ্বারা ভাগ করে ঘাতগুলোকে ছোট করি:
    \begin{align*}
        334 &= 3 \times 111 + 1 \Rightarrow \omega^{334} = \omega^1 = \omega \\
        365 &= 3 \times 121 + 2 \Rightarrow \omega^{365} = \omega^2
    \end{align*}
    সমীকরণে মান বসিয়ে পাই:
    \[ P = 4 + 5\omega + 3\omega^2 = 1 + 2\omega + 3(1 + \omega + \omega^2) \]
    আমরা জানি, $1 + \omega + \omega^2 = 0$। অতএব:
    \[ P = 1 + 2\omega + 3(0) = 1 + 2\omega \]
    এখন $\omega$ এর মূল মানটি বসিয়ে পাই:
    \[ P = 1 + 2\left(-\frac{1}{2} + i\frac{\sqrt{3}}{2}\right) = 1 - 1 + i\sqrt{3} = i\sqrt{3} \]
\end{tcolorbox}

\vspace{1em}

% ------------------------------------------
% QUESTION SECTION
% ------------------------------------------
\begin{tcolorbox}[questionbox={প্রশ্ন (Question) 24}]
    \noindent
    যদি পরিবর্তনশীল জটিল সংখ্যা $z = \frac{(1+i\sqrt{3})(\cos\theta - i\sin\theta)}{(1+i)(\cos\theta + i\sin\theta)}$ হয়, তবে $z$ এর মুখ্য আর্গুমেন্ট বা $\arg(z)$ নিচের কোনটি?
    
    \vspace{0.8em}
    \begin{english}
    \noindent\textit{\color{darkgray}
    If the variable complex number is $z = \frac{(1+i\sqrt{3})(\cos\theta - i\sin\theta)}{(1+i)(\cos\theta + i\sin\theta)}$, then the principal argument $\arg(z)$ of $z$ is:
    }
    \end{english}
\end{tcolorbox}

% ------------------------------------------
% OPTIONS SECTION
% ------------------------------------------
\begin{itemize}[label={}, leftmargin=1cm, itemsep=0.5em]
    \item[\textbf{A)}] বা $\frac{\pi}{12} - 2\theta$
    \item[\textbf{B)}] বা $\frac{\pi}{12} + 2\theta$
    \item[\textbf{C)}] বা বা $\frac{\pi}{4} - \theta$
    \item[\textbf{D)}] বা বা $\frac{\pi}{3} - 2\theta$
\end{itemize}

% ------------------------------------------
% ANSWER HIGHLIGHT
% ------------------------------------------
\vspace{0.5em}
\begin{tcolorbox}[answerbox]
    \textbf{\color{success} উত্তর:} A) $\frac{\pi}{12} - 2\theta$
\end{tcolorbox}
\vspace{1em}

% ------------------------------------------
% SOLUTION SECTION
% ------------------------------------------
\begin{tcolorbox}[solutionbox={সমাধান (Solution)}]
    \noindent
    লব ও হরের জটিল সংখ্যা এবং ত্রিকোণমিতিক অংশগুলোকে অয়লারের আকারে প্রকাশ করি:
    \begin{enumerate}
        \item বা $1+i\sqrt{3} = 2\left(\frac{1}{2} + i\frac{\sqrt{3}}{2}\right) = 2 e^{i\pi/3}$
        \item বা $\cos\theta - i\sin\theta = e^{-i\theta}$
        \item বা $1+i = \sqrt{2}\left(\frac{1}{\sqrt{2}} + i\frac{1}{\sqrt{2}}\right) = \sqrt{2} e^{i\pi/4}$
        \item বা $\cos\theta + i\sin\theta = e^{i\theta}$
    \end{enumerate}
    এখন সমীকরণে এই মানগুলো প্রতিস্থাপন করি:
    \begin{align*}
        z &= \frac{(2 e^{i\pi/3})(e^{-i\theta})}{(\sqrt{2} e^{i\pi/4})(e^{i\theta})} \\
        \Rightarrow z &= \sqrt{2} \cdot \frac{e^{i(\pi/3 - \theta)}}{e^{i(\pi/4 + \theta)}} \\
        \Rightarrow z &= \sqrt{2} \cdot e^{i(\pi/3 - \theta - \pi/4 - \theta)} \\
        \Rightarrow z &= \sqrt{2} \cdot e^{i(\pi/12 - 2\theta)}
    \end{align*}
    এটি অয়লারের সাধারণ আকৃতি $R e^{i\Phi}$ প্রকাশ করে, যেখানে ব্যাসার্ধ $R = \sqrt{2}$। অতএব, এর মুখ্য আর্গুমেন্ট হলো: $\arg(z) = \frac{\pi}{12} - 2\theta$।
\end{tcolorbox}

\vspace{1em}

% ------------------------------------------
% QUESTION SECTION
% ------------------------------------------
\begin{tcolorbox}[questionbox={প্রশ্ন (Question) 25}]
    \noindent
    যদি জটিল সমীকরণ বা $(\sqrt{3} + i)^{100} = 2^{99}(x + iy)$ সত্য হয় (যেখানে $x, y \in \mathbb{R}$), তবে মডুলাসের বৈশিষ্ট্য ব্যবহার করে $x^2 + y^2$ এর সঠিক মান নির্ণয় করো।
    
    \vspace{0.8em}
    \begin{english}
    \noindent\textit{\color{darkgray}
    If the complex equation $(\sqrt{3} + i)^{100} = 2^{99}(x + iy)$ is true (where $x, y \in \mathbb{R}$), then using the properties of modulus, determine the correct value of $x^2 + y^2$.
    }
    \end{english}
\end{tcolorbox}

% ------------------------------------------
% OPTIONS SECTION
% ------------------------------------------
\begin{itemize}[label={}, leftmargin=1cm, itemsep=0.5em]
    \item[\textbf{A)}] বা $1$
    \item[\textbf{B)}] বা $2$
    \item[\textbf{C)}] বা $4$
    \item[\textbf{D)}] বা $8$
\end{itemize}

% ------------------------------------------
% ANSWER HIGHLIGHT
% ------------------------------------------
\vspace{0.5em}
\begin{tcolorbox}[answerbox]
    \textbf{\color{success} উত্তর:} C) $4$
\end{tcolorbox}
\vspace{1em}

% ------------------------------------------
% SOLUTION SECTION
% ------------------------------------------
\begin{tcolorbox}[solutionbox={সমাধান (Solution)}]
    \noindent
    প্রদত্ত সমীকরণটি হলো: 
    \[ (\sqrt{3} + i)^{100} = 2^{99}(x + iy) \]
    উভয় পক্ষের মডুলাস বা পরমমান নিয়ে পাই:
    \[ \left|(\sqrt{3} + i)^{100}\right| = \left|2^{99}(x + iy)\right| \]
    মডুলাসের ঘাতের ধর্মাবলি প্রয়োগ করে পাই ($|Z^n| = |Z|^n$):
    \[ |\sqrt{3} + i|^{100} = 2^{99} |x + iy| \]
    আমরা জানি: 
    \[ |\sqrt{3} + i| = \sqrt{(\sqrt{3})^2 + 1^2} = \sqrt{3+1} = \sqrt{4} = 2 \]
    এবং $|x + iy| = \sqrt{x^2 + y^2}$। সমীকরণে এই মানগুলো বসিয়ে পাই:
    \[ 2^{100} = 2^{99} \sqrt{x^2 + y^2} \]
    উভয় পক্ষকে $2^{99}$ দ্বারা ভাগ করে পাই:
    \[ \sqrt{x^2 + y^2} = \frac{2^{100}}{2^{99}} = 2 \]
    উভয় পক্ষকে বর্গ করে পাই:
    \[ x^2 + y^2 = 2^2 = 4 \]
\end{tcolorbox}
% ------------------------------------------
% QUESTION SECTION
% ------------------------------------------
\begin{tcolorbox}[questionbox={প্রশ্ন (Question) 26}]
    \noindent
    যদি $a+b+c = 0$ হয়, তবে বাস্তব সহগবিশিষ্ট দ্বিঘাত সমীকরণ $3ax^2 + 2bx + c = 0$-এর মূলসমূহের প্রকৃতি ও অবস্থানের ক্ষেত্রে নিচের কোনটি সঠিক?
    
    \vspace{0.8em}
    \begin{english}
    \noindent\textit{\color{darkgray}
    If $a+b+c = 0$, then which of the following is correct regarding the nature and location of the roots of the quadratic equation $3ax^2 + 2bx + c = 0$ with real coefficients?
    }
    \end{english}
\end{tcolorbox}

% ------------------------------------------
% OPTIONS SECTION
% ------------------------------------------
\begin{itemize}[label={}, leftmargin=1cm, itemsep=0.5em]
    \item[\textbf{A)}] সমীকরণটির অন্তত একটি মূল $(0, 1)$ ব্যবধিতে থাকবে
    \item[\textbf{B)}] সমীকরণটির একটি মূল $(2, 3)$ এবং অপর মূলটি $(-2, -1)$ ব্যবধিতে থাকবে
    \item[\textbf{C)}] সমীকরণটির মূলগুলো সর্বদা কাল্পনিক বা জটিল হবে
    \item[\textbf{D)}] সমীকরণটির কোনো বাস্তব মূল থাকবে না
\end{itemize}

% ------------------------------------------
% ANSWER HIGHLIGHT
% ------------------------------------------
\vspace{0.5em}
\begin{tcolorbox}[answerbox]
    \textbf{\color{success} উত্তর:} A) সমীকরণটির অন্তত একটি মূল $(0, 1)$ ব্যবধিতে থাকবে
\end{tcolorbox}
\vspace{1em}

% ------------------------------------------
% SOLUTION SECTION
% ------------------------------------------
\begin{tcolorbox}[solutionbox={সমাধান (Solution)}]
    \noindent
    ধরি, একটি চলক অপেক্ষক $f(x) = ax^3 + bx^2 + cx$। 
    
    যেহেতু $f(x)$ একটি বহুপদী অপেক্ষক, এটি বাস্তব সংখ্যা সেটজুড়ে অবিচ্ছিন্ন এবং $(0, 1)$ ব্যবধিতে অন্তরীকরণযোগ্য। আমরা পাই:
    \begin{align*}
        f(0) &= a(0)^3 + b(0)^2 + c(0) = 0 \\
        f(1) &= a(1)^3 + b(1)^2 + c(1) = a + b + c
    \end{align*}
    প্রদত্ত শর্তানুসারে, $a+b+c = 0$, অতএব $f(1) = 0$। 
    
    যেহেতু $f(0) = f(1) = 0$, রোল-এর উপপাদ্য (Rolle's Theorem) অনুযায়ী $(0, 1)$ ব্যবধিতে অন্তত একটি বাস্তব সংখ্যা $x_0$ থাকবে যার জন্য:
    \[ f'(x_0) = 0 \]
    এখন $f(x)$-কে অন্তরীকরণ করে পাই:
    \[ f'(x) = 3ax^2 + 2bx + c \]
    অতএব, $3ax_0^2 + 2bx_0 + c = 0$ সমীকরণটির অন্তত একটি বাস্তব মূল $x_0$, যা $(0, 1)$ ব্যবধিতে অবস্থান করবে।
\end{tcolorbox}

\vspace{1em}

% ------------------------------------------
% QUESTION SECTION
% ------------------------------------------
\begin{tcolorbox}[questionbox={প্রশ্ন (Question) 27}]
    \noindent
    যদি ত্রিঘাত সমীকরণ $x^3 + 2x^2 + 10x + k = 0$-এর মূলত্রয় গুণোত্তর প্রগমনে (G.P.) থাকে, তবে $k$-এর সঠিক মান কত?
    
    \vspace{0.8em}
    \begin{english}
    \noindent\textit{\color{darkgray}
    If the roots of the cubic equation $x^3 + 2x^2 + 10x + k = 0$ are in Geometric Progression (G.P.), then what is the correct value of $k$?
    }
    \end{english}
\end{tcolorbox}

% ------------------------------------------
% OPTIONS SECTION
% ------------------------------------------
\begin{itemize}[label={}, leftmargin=1cm, itemsep=0.5em]
    \item[\textbf{A)}] $k = 125$
    \item[\textbf{B)}] $k = -125$
    \item[\textbf{C)}] $k = 27$
    \item[\textbf{D)}] $k = -27$
\end{itemize}

% ------------------------------------------
% ANSWER HIGHLIGHT
% ------------------------------------------
\vspace{0.5em}
\begin{tcolorbox}[answerbox]
    \textbf{\color{success} উত্তর:} A) $k = 125$
\end{tcolorbox}
\vspace{1em}

% ------------------------------------------
% SOLUTION SECTION
% ------------------------------------------
\begin{tcolorbox}[solutionbox={সমাধান (Solution)}]
    \noindent
    ধরি, ত্রিঘাত সমীকরণটির মূলত্রয় হলো $\frac{a}{r}$, $a$, এবং $ar$ (যেখানে $r$ গুণোত্তর প্রগমনের সাধারণ অনুপাত)। 
    
    ত্রিঘাত সমীকরণের মূল ও সহগের সম্পর্ক থেকে পাই:
    \begin{itemize}
        \item \textbf{মূলত্রয়ের গুণফল:} $\left(\frac{a}{r}\right) \cdot a \cdot (ar) = -k \implies a^3 = -k$
        \item \textbf{মূলত্রয়ের সমষ্টি:} $\frac{a}{r} + a + ar = -2 \implies a\left(\frac{1}{r} + 1 + r\right) = -2 \quad \text{--- (i)}$
        \item \textbf{দুটি করে মূলের গুণফলের সমষ্টি:} $\left(\frac{a}{r} \cdot a\right) + (a \cdot ar) + \left(ar \cdot \frac{a}{r}\right) = 10$
        \[ \implies \frac{a^2}{r} + a^2 r + a^2 = 10 \implies a^2\left(\frac{1}{r} + 1 + r\right) = 10 \quad \text{--- (ii)} \]
    \end{itemize}
    এখন সমীকরণ (ii)-কে সমীকরণ (i) দ্বারা ভাগ করে পাই:
    \[ \frac{a^2\left(\frac{1}{r} + 1 + r\right)}{a\left(\frac{1}{r} + 1 + r\right)} = \frac{10}{-2} \implies a = -5 \]
    যেহেতু $a^3 = -k$, আমরা পাই:
    \[ (-5)^3 = -k \implies -125 = -k \implies k = 125 \]
\end{tcolorbox}

\vspace{1em}

% ------------------------------------------
% QUESTION SECTION
% ------------------------------------------
\begin{tcolorbox}[questionbox={প্রশ্ন (Question) 28}]
    \noindent
    যদি $x^2 - x + 1 = 0$ সমীকরণের দুটি ভিন্ন মূল $\alpha$ ও $\beta$ হয়, তবে $\alpha^{101} + \beta^{107}$-এর মান কত হবে?
    
    \vspace{0.8em}
    \begin{english}
    \noindent\textit{\color{darkgray}
    If $\alpha$ and $\beta$ are the two distinct roots of the equation $x^2 - x + 1 = 0$, then what is the value of $\alpha^{101} + \beta^{107}$?
    }
    \end{english}
\end{tcolorbox}

% ------------------------------------------
% OPTIONS SECTION
% ------------------------------------------
\begin{itemize}[label={}, leftmargin=1cm, itemsep=0.5em]
    \item[\textbf{A)}] রাশিটির মান -1
    \item[\textbf{B)}] রাশিটির মান 0
    \item[\textbf{C)}] রাশিটির মান 1
    \item[\textbf{D)}] রাশিটির মান 2
\end{itemize}

% ------------------------------------------
% ANSWER HIGHLIGHT
% ------------------------------------------
\vspace{0.5em}
\begin{tcolorbox}[answerbox]
    \textbf{\color{success} উত্তর:} C) রাশিটির মান 1
\end{tcolorbox}
\vspace{1em}

% ------------------------------------------
% SOLUTION SECTION
% ------------------------------------------
\begin{tcolorbox}[solutionbox={সমাধান (Solution)}]
    \noindent
    প্রদত্ত সমীকরণটি হলো $x^2 - x + 1 = 0$। 
    
    আমরা জানি, এককের কাল্পনিক ঘনমূল $\omega = \frac{-1+i\sqrt{3}}{2}$ এবং $\omega^2 = \frac{-1-i\sqrt{3}}{2}$ হলে, সমীকরণ $x^2 + x + 1 = 0$-এর মূলদ্বয় হলো $\omega$ ও $\omega^2$। প্রদত্ত সমীকরণ $x^2 - x + 1 = 0$-এর মূলদ্বয় হবে:
    \[ x = -\omega \quad \text{এবং} \quad x = -\omega^2 \]
    এখন, $\alpha = -\omega$ এবং $\beta = -\omega^2$। 
    
    এখন, $\alpha^{101} + \beta^{107}$-এর মান বের করি:
    \[ \alpha^{101} + \beta^{107} = (-\omega)^{101} + (-\omega^2)^{107} = -\omega^{101} - \omega^{214} \]
    যেহেতু $\omega^3 = 1$:
    \begin{align*}
        \omega^{101} &= \omega^{99} \cdot \omega^2 = (\omega^3)^{33} \cdot \omega^2 = \omega^2 \\
        \omega^{214} &= \omega^{213} \cdot \omega = (\omega^3)^{71} \cdot \omega = \omega
    \end{align*}
    অতএব:
    \[ -\omega^{101} - \omega^{214} = -(\omega^2 + \omega) \]
    আমরা জানি, $1 + \omega + \omega^2 = 0 \implies \omega^2 + \omega = -1$। 
    
    অতএব, $-\omega^{101} - \omega^{214} = -(-1) = 1$।
\end{tcolorbox}
% ------------------------------------------
% QUESTION SECTION
% ------------------------------------------
\begin{tcolorbox}[questionbox={প্রশ্ন (Question) 29}]
    \noindent
    $c$-এর কোন বাস্তব ব্যবধির জন্য $f(x) = \frac{x^2 + 2x + c}{x^2 + 4x + 3c}$ রাশিটি সকল বাস্তব মান গ্রহণ করতে পারে?
    
    \vspace{0.8em}
    \begin{english}
    \noindent\textit{\color{darkgray}
    For what real interval of $c$ can the expression $f(x) = \frac{x^2 + 2x + c}{x^2 + 4x + 3c}$ take all real values?
    }
    \end{english}
\end{tcolorbox}

% ------------------------------------------
% OPTIONS SECTION
% ------------------------------------------
\begin{itemize}[label={}, leftmargin=1cm, itemsep=0.5em]
    \item[\textbf{A)}] $c \in (0, 1)$ ব্যবধিতে থাকতে হবে
    \item[\textbf{B)}] $c \in [0, 1]$ ব্যবধিতে থাকতে হবে
    \item[\textbf{C)}] $c \in (1, 4/3)$ ব্যবধিতে থাকতে হবে
    \item[\textbf{D)}] $c \in (-\infty, 0)$ ব্যবধিতে থাকতে হবে
\end{itemize}

% ------------------------------------------
% ANSWER HIGHLIGHT
% ------------------------------------------
\vspace{0.5em}
\begin{tcolorbox}[answerbox]
    \textbf{\color{success} উত্তর:} A) $c \in (0, 1)$ ব্যবধিতে থাকতে হবে
\end{tcolorbox}
\vspace{1em}

% ------------------------------------------
% SOLUTION SECTION
% ------------------------------------------
\begin{tcolorbox}[solutionbox={সমাধান (Solution)}]
    \noindent
    ধরি, $y = \frac{x^2 + 2x + c}{x^2 + 4x + 3c}$।
    
    বজ্রগুণন করে পাই:
    \[ y(x^2 + 4x + 3c) = x^2 + 2x + c \implies (y-1)x^2 + 2(2y-1)x + c(3y-1) = 0 \]
    যেহেতু $x$ বাস্তব, তাই এই দ্বিঘাত সমীকরণটির নিশ্চয়ক $D \ge 0$ হতে হবে সকল বাস্তব $y$-এর জন্য:
    \begin{align*}
        D &= [2(2y-1)]^2 - 4(y-1)c(3y-1) \ge 0 \\
        &\implies 4(4y^2 - 4y + 1) - 4c(3y^2 - 4y + 1) \ge 0
    \end{align*}
    উভয়পক্ষকে $4$ দ্বারা ভাগ করে পাই:
    \[ (4 - 3c)y^2 - 4(1-c)y + (1-c) \ge 0 \]
    এই অসমতাটি সকল বাস্তব $y$-এর জন্য সত্য হতে হলে দুটি শর্ত প্রয়োজন:
    \begin{enumerate}
        \item $y^2$ এর সহগ ধনাত্মক হতে হবে: $4 - 3c > 0 \implies c < \frac{4}{3}$
        \item এই $y$-এর দ্বিঘাত রাশির নিশ্চয়ক $D_y \le 0$ হতে হবে:
        \begin{align*}
            D_y &= [-4(1-c)]^2 - 4(4-3c)(1-c) \le 0 \\
            &\implies 16(1-c)^2 - 4(4-3c)(1-c) \le 0
        \end{align*}
    \end{enumerate}
    উভয়পক্ষকে $4(1-c)$ দ্বারা ভাগ করি (যেখানে $c \ne 1$):
    \[ 4(1-c) - (4-3c) \le 0 \implies 4 - 4c - 4 + 3c \le 0 \implies -c \le 0 \implies c \ge 0 \]
    অতএব, আমরা পাই $0 \le c \le 1$। 
    
    তবে $c=0$ হলে ফাংশনটি সরল হয়ে যায় এবং সকল বাস্তব মান নিতে পারে না। $c=1$ হলেও এটি সরলীকৃত হয়ে যায়। অতএব, সকল বাস্তব মান গ্রহণের জন্য সঠিক ব্যবধি হলো $c \in (0, 1)$।
\end{tcolorbox}

\vspace{1em}

% ------------------------------------------
% QUESTION SECTION
% ------------------------------------------
\begin{tcolorbox}[questionbox={প্রশ্ন (Question) 30}]
    \noindent
    $\alpha, \beta, \gamma$ সমীকরণ $x^3 - 3x + 1 = 0$-এর তিনটি ভিন্ন বাস্তব মূল হলে, $(\alpha - \beta)^2(\beta - \gamma)^2(\gamma - \alpha)^2$-এর মান কত?
    
    \vspace{0.8em}
    \begin{english}
    \noindent\textit{\color{darkgray}
    If $\alpha, \beta, \gamma$ are the three distinct real roots of the equation $x^3 - 3x + 1 = 0$, then what is the value of $(\alpha - \beta)^2(\beta - \gamma)^2(\gamma - \alpha)^2$?
    }
    \end{english}
\end{tcolorbox}

% ------------------------------------------
% OPTIONS SECTION
% ------------------------------------------
\begin{itemize}[label={}, leftmargin=1cm, itemsep=0.5em]
    \item[\textbf{A)}] রাশিটির মান $81$ হবে
    \item[\textbf{B)}] রাশিটির মান $27$ হবে
    \item[\textbf{C)}] রাশিটির মান $108$ হবে
    \item[\textbf{D)}] রাশিটির মান $54$ হবে
\end{itemize}

% ------------------------------------------
% ANSWER HIGHLIGHT
% ------------------------------------------
\vspace{0.5em}
\begin{tcolorbox}[answerbox]
    \textbf{\color{success} উত্তর:} A) রাশিটির মান $81$ হবে
\end{tcolorbox}
\vspace{1em}

% ------------------------------------------
% SOLUTION SECTION
% ------------------------------------------
\begin{tcolorbox}[solutionbox={সমাধান (Solution)}]
    \noindent
    আমরা জানি, যেকোনো সাধারণ ত্রিঘাত সমীকরণ $x^3 + px + q = 0$-এর নিশ্চায়ক (Discriminant) $\Delta$ এবং এর মূলদ্বয়ের পার্থক্যের বর্গের গুণফলের সম্পর্কটি হলো:
    \[ (\alpha - \beta)^2(\beta - \gamma)^2(\gamma - \alpha)^2 = -4p^3 - 27q^2 \]
    প্রদত্ত সমীকরণ $x^3 - 3x + 1 = 0$ এর জন্য সহগগুলো হলো:
    \[ p = -3 \quad \text{এবং} \quad q = 1 \]
    এই মানগুলো সূত্রে বসিয়ে পাই:
    \[ -4(-3)^3 - 27(1)^2 = -4(-27) - 27 = 108 - 27 = 81 \]
    অতএব, নির্ণেয় মান হলো $81$।
\end{tcolorbox}
% ------------------------------------------
% QUESTION SECTION
% ------------------------------------------
\begin{tcolorbox}[questionbox={প্রশ্ন (Question) 31}]
    \noindent
    যদি $P(x)$ একটি $n$ ($n \ge 2$) ঘাতের বাস্তব সহগবিশিষ্ট বহুপদী হয় যার প্রতিটি মূল বাস্তব ও ভিন্ন, তবে $f(x) = [P'(x)]^2 - P(x)P''(x) = 0$ সমীকরণটির কতটি বাস্তব সমাধান রয়েছে?
    
    \vspace{0.8em}
    \begin{english}
    \noindent\textit{\color{darkgray}
    If $P(x)$ is a polynomial of degree $n$ ($n \ge 2$) with real coefficients such that all its roots are real and distinct, then how many real solutions does the equation $f(x) = [P'(x)]^2 - P(x)P''(x) = 0$ have?
    }
    \end{english}
\end{tcolorbox}

% ------------------------------------------
% OPTIONS SECTION
% ------------------------------------------
\begin{itemize}[label={}, leftmargin=1cm, itemsep=0.5em]
    \item[\textbf{A)}] বাস্তব সমাধান সংখ্যা $0$ (কোনো বাস্তব সমাধান নেই)
    \item[\textbf{B)}] বাস্তব সমাধান সংখ্যা $n$
    \item[\textbf{C)}] বাস্তব সমাধান সংখ্যা $n-1$
    \item[\textbf{D)}] বাস্তব সমাধান সংখ্যা $2$
\end{itemize}

% ------------------------------------------
% ANSWER HIGHLIGHT
% ------------------------------------------
\vspace{0.5em}
\begin{tcolorbox}[answerbox]
    \textbf{\color{success} উত্তর:} A) বাস্তব সমাধান সংখ্যা $0$ (কোনো বাস্তব সমাধান নেই)
\end{tcolorbox}
\vspace{1em}

% ------------------------------------------
% SOLUTION SECTION
% ------------------------------------------
\begin{tcolorbox}[solutionbox={সমাধান (Solution)}]
    \noindent
    ধরি, $P(x)$ এর বাস্তব ও ভিন্ন মূলগুলো হলো $\alpha_1, \alpha_2, \dots, \alpha_n$। 
    
    তাহলে আমরা লিখতে পারি:
    \[ P(x) = c(x-\alpha_1)(x-\alpha_2)\dots(x-\alpha_n) \quad (\text{যেখানে } c \ne 0) \]
    উভয়পক্ষে প্রাকৃতিক লগারিদম ($\ln$) নিয়ে অন্তরীকরণ করে পাই:
    \begin{align*}
        \ln|P(x)| &= \ln|c| + \sum_{i=1}^n \ln|x-\alpha_i| \\
        \frac{P'(x)}{P(x)} &= \sum_{i=1}^n \frac{1}{x-\alpha_i}
    \end{align*}
    এখন উভয়পক্ষকে পুনরায় $x$ এর সাপেক্ষে অন্তরীকরণ করি:
    \[ \frac{P''(x)P(x) - [P'(x)]^2}{[P(x)]^2} = -\sum_{i=1}^n \frac{1}{(x-\alpha_i)^2} \]
    উভয়পক্ষকে $-1$ দ্বারা গুণ করে পাই:
    \[ \frac{[P'(x)]^2 - P(x)P''(x)}{[P(x)]^2} = \sum_{i=1}^n \frac{1}{(x-\alpha_i)^2} \]
    যেহেতু $x$ বাস্তব এবং $\alpha_i$ বাস্তব ও ভিন্ন, তাই ডানপক্ষের প্রতিটি পদ $\frac{1}{(x-\alpha_i)^2} > 0$ হবে (যেকোনো বাস্তব $x \ne \alpha_i$ এর জন্য)। 
    
    অতএব:
    \[ [P'(x)]^2 - P(x)P''(x) = [P(x)]^2 \sum_{i=1}^n \frac{1}{(x-\alpha_i)^2} > 0 \]
    আর যদি $x = \alpha_i$ (অর্থাৎ সচিব $P(x) = 0$) হয়, তবে:
    \[ [P'(x)]^2 - P(x)P''(x) = [P'(\alpha_i)]^2 \]
    যেহেতু মূলগুলো ভিন্ন, তাই সাধারণ মূলের শর্তানুসারে $P'(\alpha_i) \ne 0 \implies [P'(\alpha_i)]^2 > 0$। 
    
    সুতরাং সকল বাস্তব $x$ এর জন্য $[P'(x)]^2 - P(x)P''(x) > 0$। 
    
    অতএব, $f(x) = 0$ সমীকরণটির কোনো বাস্তব সমাধান নেই (বাস্তব সমাধান সংখ্যা $0$)।
\end{tcolorbox}
% ------------------------------------------
% QUESTION SECTION
% ------------------------------------------
\begin{tcolorbox}[questionbox={প্রশ্ন (Question) 32}]
    \noindent
    যদি সকল বাস্তব $x$-এর জন্য $(a-1)x^2 - 2(a+1)x + 3(a-1) > 0$ অসমতাটি সত্য হয়, তবে $a$-এর সঠিক ব্যবধি কোনটি?
    
    \vspace{0.8em}
    \begin{english}
    \noindent\textit{\color{darkgray}
    If the inequality $(a-1)x^2 - 2(a+1)x + 3(a-1) > 0$ is true for all real values of $x$, which is the correct interval for $a$?
    }
    \end{english}
\end{tcolorbox}

% ------------------------------------------
% OPTIONS SECTION
% ------------------------------------------
\begin{itemize}[label={}, leftmargin=1cm, itemsep=0.5em]
    \item[\textbf{A)}] $a \in (2+\sqrt{3}, \infty)$ ব্যবধিতে থাকতে হবে
    \item[\textbf{B)}] $a \in (1, 2+\sqrt{3})$ ব্যবধিতে থাকতে হবে
    \item[\textbf{C)}] $a \in (-\infty, 2-\sqrt{3})$ ব্যবধিতে থাকতে হবে
    \item[\textbf{D)}] $a \in (2-\sqrt{3}, 2+\sqrt{3})$ ব্যবধিতে থাকতে হবে
\end{itemize}

% ------------------------------------------
% ANSWER HIGHLIGHT
% ------------------------------------------
\vspace{0.5em}
\begin{tcolorbox}[answerbox]
    \textbf{\color{success} উত্তর:} A) $a \in (2+\sqrt{3}, \infty)$ ব্যবধিতে থাকতে হবে
\end{tcolorbox}
\vspace{1em}

% ------------------------------------------
% SOLUTION SECTION
% ------------------------------------------
\begin{tcolorbox}[solutionbox={সমাধান (Solution)}]
    \noindent
    একটি দ্বিঘাত রাশি $Ax^2 + Bx + C > 0$ সকল বাস্তব $x$ এর জন্য সত্য হতে হলে দুটি শর্ত প্রয়োজন:
    \begin{enumerate}
        \item $A > 0 \implies a-1 > 0 \implies a > 1$।
        \item নিশ্চয়ক $D < 0$:
        \begin{align*}
            D &= [-2(a+1)]^2 - 4(a-1)(3(a-1)) < 0 \\
            &\implies 4(a+1)^2 - 12(a-1)^2 < 0
        \end{align*}
    \end{enumerate}
    উভয়পক্ষকে $4$ দ্বারা ভাগ করে পাই:
    \begin{align*}
        (a+1)^2 - 3(a-1)^2 &< 0 \\
        &\implies a^2 + 2a + 1 - 3(a^2 - 2a + 1) < 0 \\
        &\implies a^2 + 2a + 1 - 3a^2 + 6a - 3 < 0 \\
        &\implies -2a^2 + 8a - 2 < 0
    \end{align*}
    উভয়পক্ষকে $-2$ দ্বারা ভাগ করে পাই (অসমতার দিক পরিবর্তিত হবে):
    \[ a^2 - 4a + 1 > 0 \]
    সমীকরণ $a^2-4a+1=0$ এর সমাধান হলো:
    \[ a = \frac{4 \pm \sqrt{16-4}}{2} = 2 \pm \sqrt{3} \]
    অতএব, অসমতা $a^2-4a+1 > 0$ এর সমাধান সেট হলো: $a < 2-\sqrt{3}$ অথবা $a > 2+\sqrt{3}$।
    
    প্রথম শর্ত $a > 1$ এর সাথে সমাধান সেটের ছেদবিন্দু (intersection) নির্ণয় করি: যেহেতু $2-\sqrt{3} \approx 0.268 < 1$ এবং $2+\sqrt{3} \approx 3.732 > 1$, তাই সাধারণ ব্যবধি হবে:
    \[ a > 2+\sqrt{3} \text{ বা } a \in (2+\sqrt{3}, \infty) \]
\end{tcolorbox}

\vspace{1em}

% ------------------------------------------
% QUESTION SECTION
% ------------------------------------------
\begin{tcolorbox}[questionbox={প্রশ্ন (Question) 33}]
    \noindent
    যদি $x^3 - 3x + 2 = 0$ এবং $x^4 - ax^2 + b = 0$ সমীকরণ দুটির দুটি সাধারণ মূল থাকে, তবে $(a, b)$-এর মান কত?
    
    \vspace{0.8em}
    \begin{english}
    \noindent\textit{\color{darkgray}
    If the equations $x^3 - 3x + 2 = 0$ and $x^4 - ax^2 + b = 0$ have two common roots, then what is the value of $(a, b)$?
    }
    \end{english}
\end{tcolorbox}

% ------------------------------------------
% OPTIONS SECTION
% ------------------------------------------
\begin{itemize}[label={}, leftmargin=1cm, itemsep=0.5em]
    \item[\textbf{A)}] মানটি $(5, 4)$ হবে
    \item[\textbf{B)}] মানটি $(3, 2)$ হবে
    \item[\textbf{C)}] মানটি $(4, 3)$ হবে
    \item[\textbf{D)}] মানটি $(6, 5)$ হবে
\end{itemize}

% ------------------------------------------
% ANSWER HIGHLIGHT
% ------------------------------------------
\vspace{0.5em}
\begin{tcolorbox}[answerbox]
    \textbf{\color{success} উত্তর:} A) মানটি $(5, 4)$ হবে
\end{tcolorbox}
\vspace{1em}

% ------------------------------------------
% SOLUTION SECTION
% ------------------------------------------
\begin{tcolorbox}[solutionbox={সমাধান (Solution)}]
    \noindent
    প্রথম সমীকরণটি সমাধান করি:
    \[ x^3 - 3x + 2 = 0 \implies (x-1)^2(x+2) = 0 \]
    অতএব, এর মূলগুলো হলো $x = 1, 1, -2$। 
    
    যেহেতু দ্বিতীয় সমীকরণের সাথে দুটি সাধারণ মূল রয়েছে, তাই সাধারণ মূল দুটি অবশ্যই $x = 1$ এবং $x = -2$ হতে হবে। এই মূল দুটি দ্বিতীয় সমীকরণ $x^4 - ax^2 + b = 0$ কে সিদ্ধ করবে।
    
    $x=1$ বসিয়ে পাই:
    \[ 1^4 - a(1)^2 + b = 0 \implies a - b = 1 \quad \text{--- (i)} \]
    $x=-2$ বসিয়ে পাই:
    \[ (-2)^4 - a(-2)^2 + b = 0 \implies 16 - 4a + b = 0 \implies 4a - b = 16 \quad \text{--- (ii)} \]
    এখন সমীকরণ (ii) থেকে সমীকরণ (i) বিয়োগ করি:
    \[ (4a-b) - (a-b) = 16 - 1 \implies 3a = 15 \implies a = 5 \]
    $a$-এর মান সমীকরণ (i) এ বসিয়ে পাই:
    \[ 5 - b = 1 \implies b = 4 \]
    সুতরাং, $(a, b) = (5, 4)$।
\end{tcolorbox}

\vspace{1em}

% ------------------------------------------
% QUESTION SECTION
% ------------------------------------------
\begin{tcolorbox}[questionbox={প্রশ্ন (Question) 34}]
    \noindent
    বাস্তব সহগবিশিষ্ট একটি অ-ধ্রুবক বহুপদী $P(x)$ সকল বাস্তব $x$-এর জন্য $P(x)P(x+1) = P(x^2+x+1)$ সমীকরণটি সিদ্ধ করলে, $P(x)$-এর সর্বনিম্ন ঘাত (degree) কত হতে পারে?
    
    \vspace{0.8em}
    \begin{english}
    \noindent\textit{\color{darkgray}
    If a non-constant polynomial $P(x)$ with real coefficients satisfies the equation $P(x)P(x+1) = P(x^2+x+1)$ for all real $x$, then what is the minimum possible degree of $P(x)$?
    }
    \end{english}
\end{tcolorbox}

% ------------------------------------------
% OPTIONS SECTION
% ------------------------------------------
\begin{itemize}[label={}, leftmargin=1cm, itemsep=0.5em]
    \item[\textbf{A)}] সর্বনিম্ন ঘাত $1$ হতে পারে
    \item[\textbf{B)}] সর্বনিম্ন ঘাত $2$ হতে পারে
    \item[\textbf{C)}] সর্বনিম্ন ঘাত $3$ হতে পারে
    \item[\textbf{D)}] সর্বনিম্ন ঘাত $4$ হতে পারে
\end{itemize}

% ------------------------------------------
% ANSWER HIGHLIGHT
% ------------------------------------------
\vspace{0.5em}
\begin{tcolorbox}[answerbox]
    \textbf{\color{success} উত্তর:} B) সর্বনিম্ন ঘাত $2$ হতে পারে
\end{tcolorbox}
\vspace{1em}

% ------------------------------------------
% SOLUTION SECTION
% ------------------------------------------
\begin{tcolorbox}[solutionbox={সমাধান (Solution)}]
    \noindent
    ধরি, $P(x)$ এর ঘাত $n$। তাহলে সমীকরণের উভয়পক্ষের ঘাত তুলনা করি:
    \begin{itemize}
        \item বামপক্ষের ঘাত $= \text{deg}(P(x)) + \text{deg}(P(x+1)) = n + n = 2n$।
        \item ডানপক্ষের ঘাত $= \text{deg}(P(x^2+x+1)) = n \times 2 = 2n$।
    \end{itemize}
    যেহেতু উভয়পক্ষের ঘাত সমান, তাই যেকোনো ধনাত্মক পূর্ণসংখ্যা $n \ge 1$ এর জন্য এটি সম্ভব।
    
    এখন আমরা পরীক্ষা করি $n=1$ ঘাতের জন্য কোনো বাস্তব বহুপদী পাওয়া যায় কিনা: ধরি, $P(x) = ax+b$ (যেখানে সচিব $a \ne 0$)। সমীকরণে বসিয়ে পাই:
    \begin{align*}
        (ax+b)(a(x+1)+b) &= a(x^2+x+1)+b \\
        \implies a^2x^2 + a(a+2b)x + b(a+b) &= ax^2 + ax + (a+b)
    \end{align*}
    উভয়পক্ষের সহগ তুলনা করি:
    \begin{itemize}
        \item $x^2$ এর সহগ: $a^2 = a \implies a = 1$ (যেহেতু $a \ne 0$)
        \item $x$ এর সহগ: $a(a+2b) = a \implies 1(1+2b) = 1 \implies b = 0$
        \item ধ্রুবক পদ: $b(a+b) = a+b \implies 0(1+0) = 1+0 \implies 0 = 1$ (যা অসম্ভব)
    \end{itemize}
    সুতরাং ঘাত $1$ হতে পারে না।
    
    এবার $n=2$ ঘাতের জন্য পরীক্ষা করি: ধরি, $P(x) = x^2+1$।
    \begin{itemize}
        \item \textbf{বামপক্ষ:} $P(x)P(x+1) = (x^2+1)((x+1)^2+1) = (x^2+1)(x^2+2x+2) = x^4 + 2x^3 + 3x^2 + 2x + 2$।
        \item \textbf{ডানপক্ষ:} $P(x^2+x+1) = (x^2+x+1)^2 + 1 = (x^4 + x^2 + 1 + 2x^3 + 2x + 2x^2) + 1 = x^4 + 2x^3 + 3x^2 + 2x + 2$।
    \end{itemize}
    দেখা যাচ্ছে, বামপক্ষ এবং ডানপক্ষ সমান। অতএব, $P(x) = x^2+1$ সমীকরণটি সিদ্ধ করে এবং এর ঘাত $2$। সুতরাং সর্বনিম্ন ঘাত $2$ হতে পারে।
\end{tcolorbox}
% ------------------------------------------
% QUESTION SECTION
% ------------------------------------------
\begin{tcolorbox}[questionbox={প্রশ্ন (Question) 35}]
    \noindent
    যদি $\alpha$ এবং $\beta$ সমীকরণ $x^2 + 3^{\frac{1}{4}}x + 3^{\frac{1}{2}} = 0$-এর দুটি ভিন্ন মূল হয়, তবে $\alpha^{96}(\alpha^{12}-1) + \beta^{96}(\beta^{12}-1)$ রাশিটির মান কত?
    
    \vspace{0.8em}
    \begin{english}
    \noindent\textit{\color{darkgray}
    If $\alpha$ and $\beta$ are the distinct roots of the equation $x^2 + 3^{\frac{1}{4}}x + 3^{\frac{1}{2}} = 0$, then what is the value of $\alpha^{96}(\alpha^{12}-1) + \beta^{96}(\beta^{12}-1)$?
    }
    \end{english}
\end{tcolorbox}

% ------------------------------------------
% OPTIONS SECTION
% ------------------------------------------
\begin{itemize}[label={}, leftmargin=1cm, itemsep=0.5em]
    \item[\textbf{A)}] $28 \times 3^{25}$
    \item[\textbf{B)}] $56 \times 3^{25}$
    \item[\textbf{C)}] $56 \times 3^{24}$
    \item[\textbf{D)}] $52 \times 3^{24}$
\end{itemize}

% ------------------------------------------
% ANSWER HIGHLIGHT
% ------------------------------------------
\vspace{0.5em}
\begin{tcolorbox}[answerbox]
    \textbf{\color{success} উত্তর:} D) $52 \times 3^{24}$
\end{tcolorbox}
\vspace{1em}

% ------------------------------------------
% SOLUTION SECTION
% ------------------------------------------
\begin{tcolorbox}[solutionbox={সমাধান (Solution)}]
    \noindent
    প্রদত্ত দ্বিঘাত সমীকরণটি হলো:
    \[ x^2 + 3^{\frac{1}{4}}x + \sqrt{3} = 0 \implies x^2 + \sqrt{3} = -3^{\frac{1}{4}}x \]
    উভয়পক্ষকে বর্গ করি:
    \[ (x^2 + \sqrt{3})^2 = \sqrt{3}x^2 \implies x^4 + 2\sqrt{3}x^2 + 3 = \sqrt{3}x^2 \implies x^4 + 3 = -\sqrt{3}x^2 \]
    পুনরায় উভয়পক্ষকে বর্গ করি:
    \[ (x^4 + 3)^2 = 3x^4 \implies x^8 + 6x^4 + 9 = 3x^4 \implies x^8 + 3x^4 + 9 = 0 \]
    এখন উভয়পক্ষকে $(x^4 - 3)$ দ্বারা গুণ করে পাই:
    \[ (x^4 - 3)(x^8 + 3x^4 + 9) = 0 \implies (x^4)^3 - 3^3 = 0 \implies x^{12} - 27 = 0 \implies x^{12} = 27 \]
    যেহেতু $\alpha$ এবং $\beta$ এই সমীকরণের মূল, তাই আমরা লিখতে পারি:
    \[ \alpha^{12} = 27 \quad \text{এবং} \quad \beta^{12} = 27 \]
    এখন নির্ণেয় রাশিটির মান বের করি:
    \[ E = \alpha^{96}(\alpha^{12}-1) + \beta^{96}(\beta^{12}-1) \]
    \[ \implies E = \alpha^{96}(27-1) + \beta^{96}(27-1) = 26(\alpha^{96} + \beta^{96}) \]
    যেহেতু $\alpha^{96} = (\alpha^{12})^8 = (27)^8 = (3^3)^8 = 3^{24}$ এবং $\beta^{96} = 3^{24}$, সেহেতু:
    \[ E = 26(3^{24} + 3^{24}) = 26 \times 2 \times 3^{24} = 52 \times 3^{24} \]
\end{tcolorbox}

\vspace{1em}

% ------------------------------------------
% QUESTION SECTION
% ------------------------------------------
\begin{tcolorbox}[questionbox={প্রশ্ন (Question) 36}]
    \noindent
    ধরি, $\alpha$ এবং $\beta$ সমীকরণ $x^2 - \sqrt{2}x + \sqrt{6} = 0$-এর মূলদ্বয় এবং $\frac{1}{\alpha^2}+1$ ও $\frac{1}{\beta^2}+1$ সমীকরণ $x^2 + ax + b = 0$-এর মূলদ্বয়। তবে $x^2 - (a+b-2)x + (a+b+2) = 0$ সমীকরণের মূলগুলোর প্রকৃতি কেমন হবে?
    
    \vspace{0.8em}
    \begin{english}
    \noindent\textit{\color{darkgray}
    Let $\alpha, \beta$ be the roots of the equation $x^2 - \sqrt{2}x + \sqrt{6} = 0$ and $\frac{1}{\alpha^2}+1, \frac{1}{\beta^2}+1$ be the roots of the equation $x^2 + ax + b = 0$. Then the roots of the equation $x^2 - (a+b-2)x + (a+b+2) = 0$ are:
    }
    \end{english}
\end{tcolorbox}

% ------------------------------------------
% OPTIONS SECTION
% ------------------------------------------
\begin{itemize}[label={}, leftmargin=1cm, itemsep=0.5em]
    \item[\textbf{A)}] অবাস্তব জটিল সংখ্যা
    \item[\textbf{B)}] বাস্তব এবং উভয়ই ঋণাত্মক
    \item[\textbf{C)}] বাস্তব এবং উভয়ই ধনাত্মক
    \item[\textbf{D)}] বাস্তব এবং তাদের মধ্যে ঠিক একটি ধনাত্মক
\end{itemize}

% ------------------------------------------
% ANSWER HIGHLIGHT
% ------------------------------------------
\vspace{0.5em}
\begin{tcolorbox}[answerbox]
    \textbf{\color{success} উত্তর:} B) বাস্তব এবং উভয়ই ঋণাত্মক
\end{tcolorbox}
\vspace{1em}

% ------------------------------------------
% SOLUTION SECTION
% ------------------------------------------
\begin{tcolorbox}[solutionbox={সমাধান (Solution)}]
    \noindent
    প্রদত্ত সমীকরণ $x^2 - \sqrt{2}x + \sqrt{6} = 0$ এর মূলদ্বয় $\alpha$ ও $\beta$। সহগ ও মূলের সম্পর্ক থেকে:
    \[ \alpha+\beta = \sqrt{2} \quad \text{এবং} \quad \alpha\beta = \sqrt{6} \]
    এখন, $\alpha^2+\beta^2 = (\alpha+\beta)^2 - 2\alpha\beta = 2 - 2\sqrt{6}$ এবং $\alpha^2\beta^2 = 6$।
    
    ধরি, সচিব $u = \frac{1}{\alpha^2}+1$ এবং $v = \frac{1}{\beta^2}+1$ সমীকরণ $x^2 + ax + b = 0$-এর মূলদ্বয়।
    
    মূলদ্বয়ের সমষ্টি:
    \begin{align*}
        -a &= u + v = \frac{1}{\alpha^2} + \frac{1}{\beta^2} + 2 = \frac{\alpha^2+\beta^2}{\alpha^2\beta^2} + 2 \\
        &\implies -a = \frac{2-2\sqrt{6}}{6} + 2 = \frac{1-\sqrt{6}}{3} + 2 = \frac{7-\sqrt{6}}{3} \implies a = \frac{\sqrt{6}-7}{3}
    \end{align*}
    মূলদ্বয়ের গুণফল:
    \begin{align*}
        b &= uv = \left(\frac{1}{\alpha^2}+1\right)\left(\frac{1}{\beta^2}+1\right) = \frac{1}{\alpha^2\beta^2} + \frac{\alpha^2+\beta^2}{\alpha^2\beta^2} + 1 \\
        &\implies b = \frac{1}{6} + \frac{2-2\sqrt{6}}{6} + 1 = \frac{9-2\sqrt{6}}{6}
    \end{align*}
    এখন $a+b$-এর মান নির্ণয় করি:
    \[ a+b = \frac{2\sqrt{6}-14}{6} + \frac{9-2\sqrt{6}}{6} = \frac{-5}{6} \]
    এখন নির্ণেয় সমীকরণটি হলো: $x^2 - (a+b-2)x + (a+b+2) = 0$। $a+b = -\frac{5}{6}$ বসিয়ে পাই:
    \[ x^2 - \left(-\frac{5}{6}-2\right)x + \left(-\frac{5}{6}+2\right) = 0 \implies x^2 + \frac{17}{6}x + \frac{7}{6} = 0 \implies 6x^2 + 17x + 7 = 0 \]
    এই সমীকরণটি সমাধান করি:
    \[ 6x^2 + 14x + 3x + 7 = 0 \implies 2x(3x+7) + 1(3x+7) = 0 \implies (2x+1)(3x+7) = 0 \]
    অতএব, মূলদ্বয় হলো $x = -\frac{1}{2}$ এবং $x = -\frac{7}{3}$। দেখা যাচ্ছে যে, মূলদ্বয় বাস্তব এবং উভয়ই ঋণাত্মক সংখ্যা।
\end{tcolorbox}

\vspace{1em}

% ------------------------------------------
% QUESTION SECTION
% ------------------------------------------
\begin{tcolorbox}[questionbox={প্রশ্ন (Question) 37}]
    \noindent
    ধরি, $a, b$ বাস্তব সংখ্যা এবং $a \ne 0$। যদি $ax^2 - 2bx + 5 = 0$ সমীকরণের একটি পুনরাবৃত্ত মূল (repeated root) $\alpha$ হয়, যা আবার সচিব $x^2 - 2bx - 10 = 0$ সমীকরণেরও একটি মূল হয় এবং উক্ত সমীকরণের অপর মূলটি $\beta$ হয়, তবে সচিব $\alpha^2 + \beta^2$-এর মান কত?
    
    \vspace{0.8em}
    \begin{english}
    \noindent\textit{\color{darkgray}
    Let $a, b$ be real numbers, with $a \ne 0$, such that the equation $ax^2 - 2bx + 5 = 0$ has a repeated root $\alpha$ which is also a root of the equation $x^2 - 2bx - 10 = 0$. If $\beta$ is the other root of the latter equation, then what is the value of $\alpha^2 + \beta^2$?
    }
    \end{english}
\end{tcolorbox}

% ------------------------------------------
% OPTIONS SECTION
% ------------------------------------------
\begin{itemize}[label={}, leftmargin=1cm, itemsep=0.5em]
    \item[\textbf{A)}] 25
    \item[\textbf{B)}] 24
    \item[\textbf{C)}] 26
    \item[\textbf{D)}] 28
\end{itemize}

% ------------------------------------------
% ANSWER HIGHLIGHT
% ------------------------------------------
\vspace{0.5em}
\begin{tcolorbox}[answerbox]
    \textbf{\color{success} উত্তর:} A) 25
\end{tcolorbox}
\vspace{1em}

% ------------------------------------------
% SOLUTION SECTION
% ------------------------------------------
\begin{tcolorbox}[solutionbox={সমাধান (Solution)}]
    \noindent
    যেহেতু $ax^2 - 2bx + 5 = 0$ সমীকরণের মূলদ্বয় সমান বা পুনরাবৃত্ত, তাই এর নিশ্চয়ক শূন্য হতে হবে:
    \[ D = (-2b)^2 - 4(a)(5) = 0 \implies 4b^2 - 20a = 0 \implies b^2 = 5a \]
    পুনরাবৃত্ত মূলটি হলো:
    \[ \alpha = \frac{-(-2b)}{2a} = \frac{b}{a} \]
    যেহেতু $\alpha$ আবার $x^2 - 2bx - 10 = 0$ সমীকরণেরও একটি মূল, তাই আমরা পাই:
    \[ \alpha^2 - 2b\alpha - 10 = 0 \]
    এখন $\alpha = \frac{b}{a}$ বসিয়ে পাই:
    \[ \left(\frac{b}{a}\right)^2 - 2b\left(\frac{b}{a}\right) - 10 = 0 \implies \frac{b^2}{a^2} - \frac{2b^2}{a} - 10 = 0 \]
    যেহেতু $b^2 = 5a \implies a = \frac{b^2}{5}$, এই মানটি উপরিউক্ত সমীকরণে বসিয়ে পাই:
    \[ \frac{b^2}{(b^2/5)^2} - \frac{2b^2}{b^2/5} - 10 = 0 \implies \frac{25}{b^2} - 10 - 10 = 0 \implies \frac{25}{b^2} = 20 \implies b^2 = \frac{5}{4} \]
    যেহেতু $b^2 = \frac{5}{4}$, তাই $a = \frac{b^2}{5} = \frac{1}{4}$।
    
    এখন পুনরাবৃত্ত মূল $\alpha$-কে লিখতে পারি:
    \[ \alpha = \frac{b}{a} = \frac{b}{1/4} = 4b \]
    দ্বিঘাত সমীকরণ $x^2 - 2bx - 10 = 0$ এর মূলদ্বয় $\alpha$ ও $\beta$ হওয়ায়:
    \[ \alpha + \beta = 2b \]
    যেহেতু $\alpha = 4b$, আমরা পাই:
    \[ 4b + \beta = 2b \implies \beta = -2b \]
    এখন আমাদের সচিব $\alpha^2 + \beta^2$ এর মান বের করতে হবে:
    \[ \alpha^2 + \beta^2 = (4b)^2 + (-2b)^2 = 16b^2 + 4b^2 = 20b^2 \]
    এখানে সচিব $b^2 = \frac{5}{4}$ বসিয়ে পাই:
    \[ \alpha^2 + \beta^2 = 20 \left(\frac{5}{4}\right) = 25 \]
\end{tcolorbox}
% ------------------------------------------
% QUESTION SECTION
% ------------------------------------------
\begin{tcolorbox}[questionbox={প্রশ্ন (Question) 38}]
    \noindent
    যদি $p$ এবং $q$ দুটি ধনাত্মক বাস্তব সংখ্যা হয় যেখানে $p+q = 2$ এবং $p^4 + q^4 = 272$ হয়, তবে $p$ ও $q$ নিচের কোন সমীকরণটির মূলদ্বয়?
    
    \vspace{0.8em}
    \begin{english}
    \noindent\textit{\color{darkgray}
    If $p$ and $q$ are two positive real numbers such that $p+q = 2$ and $p^4 + q^4 = 272$, then $p$ and $q$ are the roots of which of the following equations?
    }
    \end{english}
\end{tcolorbox}

% ------------------------------------------
% OPTIONS SECTION
% ------------------------------------------
\begin{itemize}[label={}, leftmargin=1cm, itemsep=0.5em]
    \item[\textbf{A)}] $x^2 - 2x + 8 = 0$
    \item[\textbf{B)}] $x^2 - 2x + 2 = 0$
    \item[\textbf{C)}] $x^2 - 2x + 16 = 0$
    \item[\textbf{D)}] $x^2 - 2x + 136 = 0$
\end{itemize}

% ------------------------------------------
% ANSWER HIGHLIGHT
% ------------------------------------------
\vspace{0.5em}
\begin{tcolorbox}[answerbox]
    \textbf{\color{success} উত্তর:} C) $x^2 - 2x + 16 = 0$
\end{tcolorbox}
\vspace{1em}

% ------------------------------------------
% SOLUTION SECTION
% ------------------------------------------
\begin{tcolorbox}[solutionbox={সমাধান (Solution)}]
    \noindent
    ধরি, নির্ণেয় সমীকরণটি হলো $x^2 - 2x + P = 0$ (যেখানে $P = pq$)। 
    
    আমরা জানি:
    \[ p^2 + q^2 = (p+q)^2 - 2pq = 2^2 - 2P = 4 - 2P \]
    পুনরায়, উচ্চতর ঘাতের সম্পর্কের সূত্রানুসারে:
    \begin{align*}
        p^4 + q^4 &= (p^2 + q^2)^2 - 2p^2q^2 \\
        &\implies p^4 + q^4 = (4 - 2P)^2 - 2P^2 = 16 - 16P + 4P^2 - 2P^2 = 2P^2 - 16P + 16
    \end{align*}
    দেওয়া আছে, $p^4 + q^4 = 272$। অতএব:
    \[ 2P^2 - 16P + 16 = 272 \implies 2P^2 - 16P - 256 = 0 \]
    উভয়পক্ষকে $2$ দ্বারা ভাগ করে পাই:
    \[ P^2 - 8P - 128 = 0 \implies (P-16)(P+8) = 0 \implies P = 16 \text{ অথবা } P = -8 \]
    যেহেতু $p$ এবং $q$ উভয়ই ধনাত্মক বাস্তব সংখ্যা, তাই এদের গুণফল $P = pq > 0$ হতে হবে। সুতরাং, $P = 16$। 
    
    নির্ণেয় সমীকরণটি হবে: $x^2 - 2x + 16 = 0$।
\end{tcolorbox}

\vspace{1em}

% ------------------------------------------
% QUESTION SECTION
% ------------------------------------------
\begin{tcolorbox}[questionbox={প্রশ্ন (Question) 39}]
    \noindent
    ধরি, $x^2 - 3x + p = 0$ সমীকরণের মূলদ্বয় $\alpha$ ও $\beta$ এবং $x^2 - 6x + q = 0$ সমীকরণের মূলদ্বয় $\gamma$ ও $\delta$। যদি $\alpha, \beta, \gamma, \delta$ একটি গুণোত্তর প্রগমন (G.P.) গঠন করে, তবে অনুপাত $(2p+q) : (2q-p)$-এর মান কত?
    
    \vspace{0.8em}
    \begin{english}
    \noindent\textit{\color{darkgray}
    Let $\alpha$ and $\beta$ be the roots of $x^2 - 3x + p = 0$ and $\gamma$ and $\delta$ be the roots of $x^2 - 6x + q = 0$. If $\alpha, \beta, \gamma, \delta$ form a geometric progression, then what is the ratio $(2p+q) : (2q-p)$?
    }
    \end{english}
\end{tcolorbox}

% ------------------------------------------
% OPTIONS SECTION
% ------------------------------------------
\begin{itemize}[label={}, leftmargin=1cm, itemsep=0.5em]
    \item[\textbf{A)}] 3:1
    \item[\textbf{B)}] 33:31
    \item[\textbf{C)}] 6:7
    \item[\textbf{D)}] 5:3
\end{itemize}

% ------------------------------------------
% ANSWER HIGHLIGHT
% ------------------------------------------
\vspace{0.5em}
\begin{tcolorbox}[answerbox]
    \textbf{\color{success} উত্তর:} C) 6:7
\end{tcolorbox}
\vspace{1em}

% ------------------------------------------
% SOLUTION SECTION
% ------------------------------------------
\begin{tcolorbox}[solutionbox={সমাধান (Solution)}]
    \noindent
    যেহেতু $\alpha, \beta$ সমীকরণ $x^2 - 3x + p = 0$-এর মূল, তাই:
    \[ \alpha+\beta = 3 \quad \text{এবং} \quad \alpha\beta = p \]
    একইভাবে, যেহেতু $\gamma, \delta$ সমীকরণ $x^2 - 6x + q = 0$-এর মূল, তাই:
    \[ \gamma+\delta = 6 \quad \text{এবং} \quad \gamma\delta = q \]
    যেহেতু $\alpha, \beta, \gamma, \delta$ গুণোত্তর প্রগমনে রয়েছে, তাই আমরা মূলগুলোকে নিম্নরূপ ধরতে পারি:
    \[ \alpha = a, \quad \beta = ar, \quad \gamma = ar^2, \quad \delta = ar^3 \quad (\text{যেখানে } r \text{ হলো সাধারণ অনুপাত}) \]
    এখন মানগুলো বসিয়ে পাই:
    \begin{align*}
        \alpha+\beta &= a + ar = a(1+r) = 3 \quad \text{--- (i)} \\
        \gamma+\delta &= ar^2 + ar^3 = ar^2(1+r) = 6 \quad \text{--- (ii)}
    \end{align*}
    সমীকরণ (ii)-কে সমীকরণ (i) দ্বারা ভাগ করে পাই:
    \[ \frac{ar^2(1+r)}{a(1+r)} = \frac{6}{3} \implies r^2 = 2 \]
    এখন $p$ ও $q$-এর মানগুলো গুণফলের মাধ্যমে প্রকাশ করি:
    \begin{align*}
        p &= \alpha\beta = a(ar) = a^2r \\
        q &= \gamma\delta = (ar^2)(ar^3) = a^2r^5
    \end{align*}
    এদের অনুপাত:
    \[ \frac{q}{p} = \frac{a^2r^5}{a^2r} = r^4 = (r^2)^2 = 2^2 = 4 \implies q = 4p \]
    আমাদের নির্ণেয় অনুপাতটি হলো: $\frac{2p+q}{2q-p}$। 
    
    এখানে $q = 4p$ বসিয়ে পাই:
    \[ \frac{2p+q}{2q-p} = \frac{2p+4p}{2(4p)-p} = \frac{6p}{8p-p} = \frac{6p}{7p} = \frac{6}{7} \]
    সুতরাং, অনুপাতটি হলো $6:7$।
\end{tcolorbox}
% ------------------------------------------
% QUESTION SECTION
% ------------------------------------------
\begin{tcolorbox}[questionbox={প্রশ্ন (Question) 40}]
    \noindent
    $m$-এর কোন বাস্তব মানসমূহের জন্য $x^2 - 2(m+1)x + m^2 + 2m - 3 = 0$ সমীকরণের উভয় মূলই $(-2, 5)$ ব্যবধিতে অবস্থান করবে?
    
    \vspace{0.8em}
    \begin{english}
    \noindent\textit{\color{darkgray}
    For what real values of $m$ do both roots of the equation $x^2 - 2(m+1)x + m^2 + 2m - 3 = 0$ lie in the interval $(-2, 5)$?
    }
    \end{english}
\end{tcolorbox}

% ------------------------------------------
% OPTIONS SECTION
% ------------------------------------------
\begin{itemize}[label={}, leftmargin=1cm, itemsep=0.5em]
    \item[\textbf{A)}] $m \in (-1, 2)$
    \item[\textbf{B)}] $m \in (-2, 5)$
    \item[\textbf{C)}] $m \in (1, 3)$
    \item[\textbf{D)}] $m \in (-2, 2)$
\end{itemize}

% ------------------------------------------
% ANSWER HIGHLIGHT
% ------------------------------------------
\vspace{0.5em}
\begin{tcolorbox}[answerbox]
    \textbf{\color{success} উত্তর:} A) $m \in (-1, 2)$
\end{tcolorbox}
\vspace{1em}

% ------------------------------------------
% SOLUTION SECTION
% ------------------------------------------
\begin{tcolorbox}[solutionbox={সমাধান (Solution)}]
    \noindent
    প্রদত্ত দ্বিঘাত সমীকরণের মূলদ্বয় বের করি:
    
    নিশ্চায়ক $D = [-2(m+1)]^2 - 4(1)(m^2+2m-3) = 4m^2 + 8m + 4 - 4m^2 - 8m + 12 = 16 > 0$। যেহেতু নিশ্চায়ক ধ্রুবক এবং ধনাত্মক, সমীকরণের মূলদ্বয় সর্বদা বাস্তব ও ভিন্ন।
    
    মূলদ্বয় হলো:
    \[ x = \frac{2(m+1) \pm \sqrt{16}}{2} = \frac{2m + 2 \pm 4}{2} = m+1 \pm 2 \]
    অতএব, মূলদ্বয় হলো $x_1 = m-1$ এবং $x_2 = m+3$। 
    
    উভয় মূল $(-2, 5)$ ব্যবধিতে থাকার শর্ত হলো:
    \[ -2 < m-1 \implies m > -1 \]
    এবং 
    \[ m+3 < 5 \implies m < 2 \]
    উভয় অসমতার সাধারণ ব্যবধি: $m \in (-1, 2)$।
\end{tcolorbox}

\vspace{1em}

% ------------------------------------------
% QUESTION SECTION
% ------------------------------------------
\begin{tcolorbox}[questionbox={প্রশ্ন (Question) 41}]
    \noindent
    যদি সচিব $x^3 - kx^2 + 36x - 36 = 0$ সমীকরণের মূলত্রয় বিপরীত প্রগমনে (H.P.) থাকে, তবে $k$-এর সঠিক বাস্তব মান কত?
    
    \vspace{0.8em}
    \begin{english}
    \noindent\textit{\color{darkgray}
    If the roots of the equation $x^3 - kx^2 + 36x - 36 = 0$ are in Harmonic Progression (H.P.), then what is the correct real value of $k$?
    }
    \end{english}
\end{tcolorbox}

% ------------------------------------------
% OPTIONS SECTION
% ------------------------------------------
\begin{itemize}[label={}, leftmargin=1cm, itemsep=0.5em]
    \item[\textbf{A)}] $k = 11$
    \item[\textbf{B)}] $k = 12$
    \item[\textbf{C)}] $k = 9$
    \item[\textbf{D)}] $k = 15$
\end{itemize}

% ------------------------------------------
% ANSWER HIGHLIGHT
% ------------------------------------------
\vspace{0.5em}
\begin{tcolorbox}[answerbox]
    \textbf{\color{success} উত্তর:} A) $k = 11$
\end{tcolorbox}
\vspace{1em}

% ------------------------------------------
% SOLUTION SECTION
% ------------------------------------------
\begin{tcolorbox}[solutionbox={সমাধান (Solution)}]
    \noindent
    যদি সমীকরণের মূলত্রয় $x_1, x_2, x_3$ বিপরীত প্রগমনে থাকে, তবে $x = \frac{1}{y}$ বসালে নতুন সমীকরণের মূলগুলো সমান্তর প্রগমনে (A.P.) থাকবে। 
    
    নতুন সমীকরণটি দাঁড়ায়:
    \[ \left(\frac{1}{y}\right)^3 - k\left(\frac{1}{y}\right)^2 + 36\left(\frac{1}{y}\right) - 36 = 0 \implies 36y^3 - 36y^2 + ky - 1 = 0 \]
    ধরি, এই সমীকরণের সমান্তর প্রগমনে থাকা মূলত্রয় হলো সচিব $a-d, a, a+d$। 
    
    মূলসমূহের সমষ্টি থেকে পাই:
    \[ 3a = \frac{36}{36} = 1 \implies a = \frac{1}{3} \]
    যেহেতু $a = 1/3$ সমীকরণটির একটি মূল, এটি সমীকরণটিকে সিদ্ধ করবে:
    \begin{align*}
        36\left(\frac{1}{3}\right)^3 - 36\left(\frac{1}{3}\right)^2 + k\left(\frac{1}{3}\right) - 1 &= 0 \\
        \implies \frac{36}{27} - \frac{36}{9} + \frac{k}{3} - 1 &= 0 \\
        \implies \frac{4}{3} - 4 + \frac{k}{3} - 1 &= 0 \implies \frac{k}{3} - \frac{11}{3} = 0 \implies k = 11
    \end{align*}
\end{tcolorbox}

\vspace{1em}

% ------------------------------------------
% QUESTION SECTION
% ------------------------------------------
\begin{tcolorbox}[questionbox={প্রশ্ন (Question) 42}]
    \noindent
    $a$-এর কোন বাস্তব মানসমূহের জন্য সচিব $x^2 - ax + a + 3 = 0$ সমীকরণের উভয় মূলই বাস্তব এবং ঋণাত্মক হবে?
    
    \vspace{0.8em}
    \begin{english}
    \noindent\textit{\color{darkgray}
    For what real values of $a$ are both roots of the equation $x^2 - ax + a + 3 = 0$ real and negative?
    }
    \end{english}
\end{tcolorbox}

% ------------------------------------------
% OPTIONS SECTION
% ------------------------------------------
\begin{itemize}[label={}, leftmargin=1cm, itemsep=0.5em]
    \item[\textbf{A)}] $a \in (-3, -2]$
    \item[\textbf{B)}] $a \in [-2, 6]$
    \item[\textbf{C)}] $a \in (-\infty, -3) \cup [6, \infty)$
    \item[\textbf{D)}] $a \in (-3, 6)$
\end{itemize}

% ------------------------------------------
% ANSWER HIGHLIGHT
% ------------------------------------------
\vspace{0.5em}
\begin{tcolorbox}[answerbox]
    \textbf{\color{success} উত্তর:} A) $a \in (-3, -2]$
\end{tcolorbox}
\vspace{1em}

% ------------------------------------------
% SOLUTION SECTION
% ------------------------------------------
\begin{tcolorbox}[solutionbox={সমাধান (Solution)}]
    \noindent
    ধরি, সমীকরণের মূলদ্বয় $\alpha$ এবং $\beta$। উভয় মূল বাস্তব এবং ঋণাত্মক হওয়ার শর্তসমূহ:
    \begin{enumerate}
        \item নিশ্চয়ক $D \ge 0$:
        \[ (-a)^2 - 4(1)(a+3) \ge 0 \implies a^2 - 4a - 12 \ge 0 \implies (a-6)(a+2) \ge 0 \]
        \[ \implies a \le -2 \text{ অথবা } a \ge 6 \]
        \item মূলদ্বয়ের সমষ্টি ঋণাত্মক হতে হবে ($\alpha+\beta < 0$):
        \[ \alpha+\beta = a < 0 \]
        \item মূলদ্বয়ের গুণফল ধনাত্মক হতে হবে ($\alpha\beta > 0$):
        \[ \alpha\beta = a+3 > 0 \implies a > -3 \]
    \end{enumerate}
    এখন শর্ত ১, ২ এবং ৩ এর সাধারণ ব্যবধি (Intersection) নির্ণয় করি:
    \begin{itemize}
        \item $a \le -2$ অথবা $a \ge 6$
        \item $a < 0$
        \item $a > -3$
    \end{itemize}
    সাধারণ অংশটি হলো: $-3 < a \le -2$ বা $a \in (-3, -2]$।
\end{tcolorbox}

\vspace{1em}

% ------------------------------------------
% QUESTION SECTION
% ------------------------------------------
\begin{tcolorbox}[questionbox={প্রশ্ন (Question) 43}]
    \noindent
    $x^2 - (k-3)x + k = 0$ সমীকরণের ঠিক একটি মূল $(1, 2)$ ব্যবধিতে থাকার জন্য বাস্তব পরামিতি $k$-এর সঠিক ব্যবধি কোনটি?
    
    \vspace{0.8em}
    \begin{english}
    \noindent\textit{\color{darkgray}
    What is the correct interval of the real parameter $k$ such that the equation $x^2 - (k-3)x + k = 0$ has exactly one root in the interval $(1, 2)$?
    }
    \end{english}
\end{tcolorbox}

% ------------------------------------------
% OPTIONS SECTION
% ------------------------------------------
\begin{itemize}[label={}, leftmargin=1cm, itemsep=0.5em]
    \item[\textbf{A)}] $k > 10$
    \item[\textbf{B)}] $k < 9$
    \item[\textbf{C)}] $9 < k < 10$
    \item[\textbf{D)}] $k > 9$
\end{itemize}

% ------------------------------------------
% ANSWER HIGHLIGHT
% ------------------------------------------
\vspace{0.5em}
\begin{tcolorbox}[answerbox]
    \textbf{\color{success} উত্তর:} A) $k > 10$
\end{tcolorbox}
\vspace{1em}

% ------------------------------------------
% SOLUTION SECTION
% ------------------------------------------
\begin{tcolorbox}[solutionbox={সমাধান (Solution)}]
    \noindent
    ধরি, $f(x) = x^2 - (k-3)x + k$। 
    
    ঠিক একটি মূল $(1, 2)$ ব্যবধিতে থাকার প্রয়োজনীয় শর্ত হলো: $f(1)f(2) < 0$। 
    
    ফাংশনের মানগুলো বসিয়ে পাই:
    \[ f(1) = 1^2 - (k-3)(1) + k = 1 - k + 3 + k = 4 > 0 \]
    যেহেতু $f(1) = 4 > 0$ (সর্বদা ধনাত্মক), তাই গুণফল ঋণাত্মক হতে হলে অবশ্যই: $f(2) < 0$ হতে হবে।
    \[ f(2) = 2^2 - (k-3)(2) + k = 4 - 2k + 6 + k = 10 - k < 0 \implies k > 10 \]
    এখন নিশ্চায়ক পরীক্ষা করি:
    \[ D = (k-3)^2 - 4k = k^2 - 10k + 9 = (k-1)(k-9) \]
    যেহেতু $k > 10$, তাই $D > 0$ সত্য, অর্থাৎ মূলগুলো বাস্তব এবং অসমান। 
    
    সুতরাং, নির্ণেয় ব্যবধি হলো $k > 10$।
\end{tcolorbox}

\vspace{1em}

% ------------------------------------------
% QUESTION SECTION
% ------------------------------------------
\begin{tcolorbox}[questionbox={প্রশ্ন (Question) 44}]
    \noindent
    যদি $x^3 - 3x^2 - 6x + k = 0$ সমীকরণের মূলত্রয় $\alpha, \beta, \gamma$ সমান্তর প্রগমনে (A.P.) থাকে, তবে $k$-এর সঠিক মান এবং মূলসমূহের ঘনফলের সমষ্টি $\alpha^3 + \beta^3 + \gamma^3$-এর মান কত?
    
    \vspace{0.8em}
    \begin{english}
    \noindent\textit{\color{darkgray}
    If the roots $\alpha, \beta, \gamma$ of the equation $x^3 - 3x^2 - 6x + k = 0$ are in Arithmetic Progression (A.P.), then what is the correct value of $k$ and the sum of the cubes of the roots $\alpha^3 + \beta^3 + \gamma^3$?
    }
    \end{english}
\end{tcolorbox}

% ------------------------------------------
% OPTIONS SECTION
% ------------------------------------------
\begin{itemize}[label={}, leftmargin=1cm, itemsep=0.5em]
    \item[\textbf{A)}] $k = 8, \alpha^3 + \beta^3 + \gamma^3 = 57$
    \item[\textbf{B)}] $k = -8, \alpha^3 + \beta^3 + \gamma^3 = 57$
    \item[\textbf{C)}] $k = 8, \alpha^3 + \beta^3 + \gamma^3 = 35$
    \item[\textbf{D)}] $k = -8, \alpha^3 + \beta^3 + \gamma^3 = 35$
\end{itemize}

% ------------------------------------------
% ANSWER HIGHLIGHT
% ------------------------------------------
\vspace{0.5em}
\begin{tcolorbox}[answerbox]
    \textbf{\color{success} উত্তর:} A) $k = 8, \alpha^3 + \beta^3 + \gamma^3 = 57$
\end{tcolorbox}
\vspace{1em}

% ------------------------------------------
% SOLUTION SECTION
% ------------------------------------------
\begin{tcolorbox}[solutionbox={সমাধান (Solution)}]
    \noindent
    ধরি, মূলত্রয় হলো $a-d, a, a+d$। সমষ্টির সম্পর্ক থেকে পাই:
    \[ 3a = 3 \implies a = 1 \]
    যেহেতু $a = 1$ সমীকরণের একটি মূল, তাই এটি সমীকরণটিকে সিদ্ধ করে:
    \[ 1^3 - 3(1)^2 - 6(1) + k = 0 \implies 1 - 3 - 6 + k = 0 \implies k = 8 \]
    সমীকরণটি দাঁড়ায়:
    \[ x^3 - 3x^2 - 6x + 8 = 0 \implies (x-1)(x^2 - 2x - 8) = 0 \implies (x-1)(x-4)(x+2) = 0 \]
    অতএব, মূল তিনটি হলো যথাক্রমে $\alpha = -2$, $\beta = 1$, $\gamma = 4$। 
    
    মূলসমূহের ঘনফলের সমষ্টি:
    \[ \alpha^3 + \beta^3 + \gamma^3 = (-2)^3 + 1^3 + 4^3 = -8 + 1 + 64 = 57 \]
\end{tcolorbox}
% ------------------------------------------
% QUESTION SECTION
% ------------------------------------------
\begin{tcolorbox}[questionbox={প্রশ্ন (Question) 45}]
    \noindent
    যদি $x^3 - 3x^2 - 9x + k = 0$ সমীকরণের মূলত্রয় সমান্তর প্রগমনে (A.P.) থাকে, তবে $k$-এর সঠিক বাস্তব মান কত?
    
    \vspace{0.8em}
    \begin{english}
    \noindent\textit{\color{darkgray}
    If the roots of the equation $x^3 - 3x^2 - 9x + k = 0$ are in Arithmetic Progression (A.P.), then what is the correct real value of $k$?
    }
    \end{english}
\end{tcolorbox}

% ------------------------------------------
% OPTIONS SECTION
% ------------------------------------------
\begin{itemize}[label={}, leftmargin=1cm, itemsep=0.5em]
    \item[\textbf{A)}] $k = 11$
    \item[\textbf{B)}] $k = -11$
    \item[\textbf{C)}] $k = 9$
    \item[\textbf{D)}] $k = -9$
\end{itemize}

% ------------------------------------------
% ANSWER HIGHLIGHT
% ------------------------------------------
\vspace{0.5em}
\begin{tcolorbox}[answerbox]
    \textbf{\color{success} উত্তর:} A) $k = 11$
\end{tcolorbox}
\vspace{1em}

% ------------------------------------------
% SOLUTION SECTION
% ------------------------------------------
\begin{tcolorbox}[solutionbox={সমাধান (Solution)}]
    \noindent
    ধরি, সমান্তর প্রগমনে থাকা মূলত্রয় হলো যথাক্রমে $a-d, a, a+d$। 
    
    সহগ ও মূলের সম্পর্ক থেকে পাই:
    \[ (a-d) + a + (a+d) = 3 \implies 3a = 3 \implies a = 1 \]
    যেহেতু $a = 1$ সমীকরণটির একটি বাস্তব মূল, তাই এটি সমীকরণটিকে সিদ্ধ করবে:
    \begin{align*}
        1^3 - 3(1)^2 - 9(1) + k &= 0 \\
        \implies 1 - 3 - 9 + k &= 0 \\
        \implies -11 + k &= 0 \implies k = 11
    \end{align*}
\end{tcolorbox}
% ------------------------------------------
% QUESTION SECTION
% ------------------------------------------
\begin{tcolorbox}[questionbox={প্রশ্ন (Question) 46}]
    \noindent
    যদি $z^2 + (3+4i)z + k = 0$ সমীকরণের একটি মূল সম্পূর্ণ বাস্তব এবং অপর মূলটি সম্পূর্ণ কাল্পনিক হয়, তবে জটিল ধ্রুবক $k$-এর সঠিক মান কোনটি?
    
    \vspace{0.8em}
    \begin{english}
    \noindent\textit{\color{darkgray}
    If one root of the equation $z^2 + (3+4i)z + k = 0$ is purely real and the other is purely imaginary, then which is the correct value of the complex constant $k$?
    }
    \end{english}
\end{tcolorbox}

% ------------------------------------------
% OPTIONS SECTION
% ------------------------------------------
\begin{itemize}[label={}, leftmargin=1cm, itemsep=0.5em]
    \item[\textbf{A)}] $k = 12i$
    \item[\textbf{B)}] $k = -12i$
    \item[\textbf{C)}] $k = 12$
    \item[\textbf{D)}] $k = -12$
\end{itemize}

% ------------------------------------------
% ANSWER HIGHLIGHT
% ------------------------------------------
\vspace{0.5em}
\begin{tcolorbox}[answerbox]
    \textbf{\color{success} উত্তর:} A) $k = 12i$
\end{tcolorbox}
\vspace{1em}

% ------------------------------------------
% SOLUTION SECTION
% ------------------------------------------
\begin{tcolorbox}[solutionbox={সমাধান (Solution)}]
    \noindent
    ধরি, সমীকরণটির মূলদ্বয় যথাক্রমে বাস্তব মূল $r \in \mathbb{R}$ এবং কাল্পনিক মূল $is \in i\mathbb{R}$ (যেখানে $s \in \mathbb{R}$)। 
    
    সহগ ও মূলের সম্পর্ক থেকে মূলদ্বয়ের সমষ্টি:
    \[ r + is = -(3+4i) = -3 - 4i \]
    উভয়পক্ষের বাস্তব ও কাল্পনিক অংশ তুলনা করে পাই: $r = -3$ এবং $s = -4$। 
    
    অতএব, সমীকরণের মূলদ্বয় হলো $-3$ এবং $-4i$। 
    
    সহগ ও মূলের সম্পর্ক থেকে মূলদ্বয়ের গুণফল:
    \[ k = (-3)(-4i) = 12i \]
\end{tcolorbox}

\vspace{1em}

% ------------------------------------------
% QUESTION SECTION
% ------------------------------------------
\begin{tcolorbox}[questionbox={প্রশ্ন (Question) 47}]
    \noindent
    $P(x) = x^{100} + x^{50} + 1$ বহুপদীটিকে $x^2 + 1$ দ্বারা ভাগ করলে ভাগশেষ কত হবে?
    
    \vspace{0.8em}
    \begin{english}
    \noindent\textit{\color{darkgray}
    What is the remainder when the polynomial $P(x) = x^{100} + x^{50} + 1$ is divided by $x^2 + 1$?
    }
    \end{english}
\end{tcolorbox}

% ------------------------------------------
% OPTIONS SECTION
% ------------------------------------------
\begin{itemize}[label={}, leftmargin=1cm, itemsep=0.5em]
    \item[\textbf{A)}] ভাগশেষ $1$
    \item[\textbf{B)}] ভাগশেষ $x+1$
    \item[\textbf{C)}] ভাগশেষ $-1$
    \item[\textbf{D)}] ভাগশেষ $0$
\end{itemize}

% ------------------------------------------
% ANSWER HIGHLIGHT
% ------------------------------------------
\vspace{0.5em}
\begin{tcolorbox}[answerbox]
    \textbf{\color{success} উত্তর:} A) ভাগশেষ $1$
\end{tcolorbox}
\vspace{1em}

% ------------------------------------------
% SOLUTION SECTION
% ------------------------------------------
\begin{tcolorbox}[solutionbox={সমাধান (Solution)}]
    \noindent
    ভাজক $x^2 + 1 = 0$ এর মূলদ্বয় হলো $\pm i$। 
    
    ভাগশেষ উপপাদ্য অনুসারে, ধরি নির্ণেয় ভাগশেষটি হলো $R(x) = ax + b$। আমরা লিখতে পারি:
    \[ x^{100} + x^{50} + 1 = Q(x)(x^2 + 1) + ax + b \]
    এখন সমীকরণটিতে $x = i$ বসিয়ে পাই:
    \[ i^{100} + i^{50} + 1 = ai + b \]
    আমরা জানি, $i^{100} = (i^4)^{25} = 1$ এবং $i^{50} = (i^2)^{25} = (-1)^{25} = -1$। অতএব:
    \[ 1 - 1 + 1 = ai + b \implies 1 = ai + b \]
    উভয়পক্ষের বাস্তব ও কাল্পনিক অংশ তুলনা করে পাই: $a = 0$ এবং সচিব $b = 1$। 
    
    সুতরাং, নির্ণেয় ভাগশেষটি হলো $1$।
\end{tcolorbox}

\vspace{1em}

% ------------------------------------------
% QUESTION SECTION
% ------------------------------------------
\begin{tcolorbox}[questionbox={প্রশ্ন (Question) 48}]
    \noindent
    যদি $x^2 - 2x\cos\theta + 1 = 0$ হয়, তবে $x^{24} + \frac{1}{x^{24}}$-এর মান নিচের কোনটি?
    
    \vspace{0.8em}
    \begin{english}
    \noindent\textit{\color{darkgray}
    If $x^2 - 2x\cos\theta + 1 = 0$, then which of the following is the value of $x^{24} + \frac{1}{x^{24}}$?
    }
    \end{english}
\end{tcolorbox}

% ------------------------------------------
% OPTIONS SECTION
% ------------------------------------------
\begin{itemize}[label={}, leftmargin=1cm, itemsep=0.5em]
    \item[\textbf{A)}] $2\cos(24\theta)$
    \item[\textbf{B)}] $\cos(24\theta)$
    \item[\textbf{C)}] $2\sin(24\theta)$
    \item[\textbf{D)}] $1$
\end{itemize}

% ------------------------------------------
% ANSWER HIGHLIGHT
% ------------------------------------------
\vspace{0.5em}
\begin{tcolorbox}[answerbox]
    \textbf{\color{success} উত্তর:} A) $2\cos(24\theta)$
\end{tcolorbox}
\vspace{1em}

% ------------------------------------------
% SOLUTION SECTION
% ------------------------------------------
\begin{tcolorbox}[solutionbox={সমাধান (Solution)}]
    \noindent
    প্রদত্ত সমীকরণ: $x^2 - 2x\cos\theta + 1 = 0$। 
    
    শ্রীধর আচার্যের সূত্র ব্যবহার করে পাই:
    \[ x = \frac{2\cos\theta \pm \sqrt{4\cos^2\theta - 4}}{2} = \cos\theta \pm i\sin\theta = e^{\pm i\theta} \]
    ধরি, $x = e^{i\theta}$। ডি ময়ভারের উপপাদ্য (De Moivre's Theorem) প্রয়োগ করে পাই:
    \[ x^{24} = (e^{i\theta})^{24} = e^{i 24\theta} \quad \text{এবং} \quad \frac{1}{x^{24}} = e^{-i 24\theta} \]
    অতএব:
    \[ x^{24} + \frac{1}{x^{24}} = e^{i 24\theta} + e^{-i 24\theta} = 2\cos(24\theta) \]
\end{tcolorbox}

\vspace{1em}

% ------------------------------------------
% QUESTION SECTION
% ------------------------------------------
\begin{tcolorbox}[questionbox={প্রশ্ন (Question) 49}]
    \noindent
    যদি $\alpha_1, \alpha_2, \dots, \alpha_{11}$ সমীকরণ $x^{11} + x^{10} + \dots + x + 1 = 0$-এর মূলসমূহ হয়, তবে $\sum_{i=1}^{11} \frac{1}{1-\alpha_i}$ রাশিটির মান কত?
    
    \vspace{0.8em}
    \begin{english}
    \noindent\textit{\color{darkgray}
    If $\alpha_1, \alpha_2, \dots, \alpha_{11}$ are the roots of the equation $x^{11} + x^{10} + \dots + x + 1 = 0$, then what is the value of $\sum_{i=1}^{11} \frac{1}{1-\alpha_i}$?
    }
    \end{english}
\end{tcolorbox}

% ------------------------------------------
% OPTIONS SECTION
% ------------------------------------------
\begin{itemize}[label={}, leftmargin=1cm, itemsep=0.5em]
    \item[\textbf{A)}] সমষ্টির মান $5.5$
    \item[\textbf{B)}] সমষ্টির মান $11$
    \item[\textbf{C)}] সমষ্টির মান $6$
    \item[\textbf{D)}] সমষ্টির মান $6.5$
\end{itemize}

% ------------------------------------------
% ANSWER HIGHLIGHT
% ------------------------------------------
\vspace{0.5em}
\begin{tcolorbox}[answerbox]
    \textbf{\color{success} উত্তর:} A) সমষ্টির মান $5.5$
\end{tcolorbox}
\vspace{1em}

% ------------------------------------------
% SOLUTION SECTION
% ------------------------------------------
\begin{tcolorbox}[solutionbox={সমাধান (Solution)}]
    \noindent
    দেওয়া আছে সমীকরণটি হলো: 
    \[ x^{11} + x^{10} + \dots + x + 1 = 0 \implies \frac{x^{12}-1}{x-1} = 0 \]
    ধরি, $y = \frac{1}{1-x} \implies 1-x = \frac{1}{y} \implies x = 1 - \frac{1}{y} = \frac{y-1}{y}$। 
    
    এই মানটি সমীকরণে বসালে আমরা পাই:
    \[ \sum_{k=0}^{11} \left(\frac{y-1}{y}\right)^k = 0 \implies \sum_{k=0}^{11} y^{11-k}(y-1)^k = 0 \]
    এখন সমীকরণটির সর্বোচ্চ ঘাত $y^{11}$ এবং তার ঠিক পরের ঘাত $y^{10}$ এর সহগ বের করি: প্রতিটি $k$ এর জন্য,
    \[ y^{11-k}(y-1)^k = y^{11-k}(y^k - ky^{k-1} + \dots) = y^{11} - ky^{10} + \dots \]
    সম্পূর্ণ সমষ্টি দাঁড়ায়:
    \[ \sum_{k=0}^{11} (y^{11} - ky^{10} + \dots) = 12y^{11} - \left(\sum_{k=0}^{11} k\right) y^{10} + \dots = 0 \]
    আমরা জানি, $\sum_{k=0}^{11} k = \frac{11 \times 12}{2} = 66$। 
    
    সমীকরণটি দাঁড়ায়: $12y^{11} - 66y^{10} + \dots = 0$। 
    
    সহগ ও মূলের সম্পর্ক থেকে নতুন মূলসমূহের সমষ্টি:
    \[ \sum_{i=1}^{11} \frac{1}{1-\alpha_i} = -\frac{-66}{12} = \frac{11}{2} = 5.5 \]
\end{tcolorbox}

\vspace{1em}

% ------------------------------------------
% QUESTION SECTION
% ------------------------------------------
\begin{tcolorbox}[questionbox={প্রশ্ন (Question) 50}]
    \noindent
    যদি $z = \cos\frac{2\pi}{7} + i\sin\frac{2\pi}{7}$ হয়, তবে $\alpha = z + z^2 + z^4$ এবং $\beta = z^3 + z^5 + z^6$ মূলদ্বয়বিশিষ্ট দ্বিঘাত সমীকরণটি নিচের কোনটি?
    
    \vspace{0.8em}
    \begin{english}
    \noindent\textit{\color{darkgray}
    If $z = \cos\frac{2\pi}{7} + i\sin\frac{2\pi}{7}$, then which of the following is the quadratic equation whose roots are $\alpha = z + z^2 + z^4$ and $\beta = z^3 + z^5 + z^6$?
    }
    \end{english}
\end{tcolorbox}

% ------------------------------------------
% OPTIONS SECTION
% ------------------------------------------
\begin{itemize}[label={}, leftmargin=1cm, itemsep=0.5em]
    \item[\textbf{A)}] $x^2 + x + 2 = 0$
    \item[\textbf{B)}] $x^2 - x + 2 = 0$
    \item[\textbf{C)}] $x^2 + x - 2 = 0$
    \item[\textbf{D)}] $x^2 + 2x + 2 = 0$
\end{itemize}

% ------------------------------------------
% ANSWER HIGHLIGHT
% ------------------------------------------
\vspace{0.5em}
\begin{tcolorbox}[answerbox]
    \textbf{\color{success} উত্তর:} A) $x^2 + x + 2 = 0$
\end{tcolorbox}
\vspace{1em}

% ------------------------------------------
% SOLUTION SECTION
% ------------------------------------------
\begin{tcolorbox}[solutionbox={সমাধান (Solution)}]
    \noindent
    যেহেতু $z = \cos\frac{2\pi}{7} + i\sin\frac{2\pi}{7} = e^{i \frac{2\pi}{7}}$, এটি এককের একটি সপ্তম কাল্পনিক মূল, তাই সচিব $z^7 = 1$ এবং $\sum_{k=1}^6 z^k = -1$। 
    
    এখন মূলদ্বয়ের সমষ্টি বের করি:
    \[ \alpha + \beta = (z + z^2 + z^4) + (z^3 + z^5 + z^6) = \sum_{k=1}^6 z^k = -1 \]
    মূলদ্বয়ের গুণফল বের করি:
    \[ \alpha\beta = (z + z^2 + z^4)(z^3 + z^5 + z^6) = z^4 + z^6 + z^7 + z^5 + z^7 + z^8 + z^7 + z^9 + z^{10} \]
    যেহেতু $z^7 = 1$, তাই আমরা পাই: $z^8 = z$, $z^9 = z^2$, $z^{10} = z^3$। অতএব:
    \[ \alpha\beta = z^4 + z^6 + 1 + z^5 + 1 + z + 1 + z^2 + z^3 = 3 + (z + z^2 + z^3 + z^4 + z^5 + z^6) \]
    \[ \implies \alpha\beta = 3 - 1 = 2 \]
    সুতরাং, দ্বিঘাত সমীকরণটি হবে:
    \[ x^2 - (\alpha+\beta)x + \alpha\beta = 0 \implies x^2 - (-1)x + 2 = 0 \implies x^2 + x + 2 = 0 \]
\end{tcolorbox}
\end{document}